\documentclass[11pt,a4paper]{article}
\usepackage[margin=23mm,headheight=15pt,headsep=8mm,footskip=11mm]{geometry}
\usepackage{amsmath,amsthm,fontspec,unicode-math}
\setmainfont{TeX Gyre Pagella}
\setsansfont{TeX Gyre Heros}
\setmonofont{Latin Modern Mono}[Scale=MatchLowercase]
\setmathfont{TeX Gyre Pagella Math}
\usepackage{microtype,booktabs,array,enumitem,xcolor,fancyhdr,needspace,hyperref}
\definecolor{inkblue}{HTML}{243F56}
\definecolor{rulegray}{HTML}{AEB7BE}
\hypersetup{colorlinks=true,linkcolor=inkblue,citecolor=inkblue,urlcolor=inkblue,
  pdftitle={Rank-four finite T-data: a computer-assisted classification},
  pdfsubject={Thirty-seven families, explicit matrices and exact certificates}}
\pagestyle{fancy}\fancyhf{}
\fancyhead[L]{\small\sffamily Rank-four finite T-data}
\fancyhead[R]{\small\sffamily Classification and certificates}
\fancyfoot[C]{\small\thepage}
\renewcommand{\headrulewidth}{0.3pt}
\setlength{\parindent}{1.2em}
\setlength{\parskip}{3pt plus 1pt}
\setlength{\emergencystretch}{1em}
\setlist{itemsep=3pt,topsep=5pt}
\raggedbottom
\newtheorem{theorem}{Theorem}[section]
\newtheorem{lemma}[theorem]{Lemma}
\newtheorem{corollary}[theorem]{Corollary}
\theoremstyle{definition}
\newtheorem{definition}[theorem]{Definition}
\newtheorem{computation}[theorem]{Exact computation}
\DeclareMathOperator{\diag}{diag}
\newcommand{\Q}{\mathbb Q}
\newcommand{\Z}{\mathbb Z}
\newcommand{\R}{\mathbb R}
\newcommand{\code}[1]{\texttt{\detokenize{#1}}}

\begin{document}
\thispagestyle{plain}
\begin{center}
{\LARGE\bfseries Rank-four finite T-data\par}
\vspace{4pt}{\large A computer-assisted classification\par}
\vspace{9pt}{\small 5 September 2026}
\end{center}
\begin{abstract}
With identity symmetrizer and diagonal leading matrix, there are thirty-seven
indecomposable rank-four families, up to rational rescaling of time, species
shifts, index permutations and sign exchange. We give explicit representatives,
their matrices at $z=1$, and exact periodicity certificates. A bound of $15$ on
off-diagonal coefficient sums reduces necessity to a finite enumeration,
followed by unbounded linear-arithmetic feasibility checks. A leading-term
obstruction excludes a remaining chain pattern. All surviving lift spaces
consist precisely of rescalings and shifts of their representatives.
\end{abstract}

\section{The classification}\label{sec:theorem}
The scope is the identity-symmetrizer, diagonal-leading-matrix setting of
Mizuno's rank-two paper \cite{Mizuno}. Throughout, the rank is the number of
species in the difference equation; the reconstructed quiver can have more
vertices.

\begin{definition}\label{def:datum}
Let $r_1,\ldots,r_4$ be positive integers and let
\[
 N_0(z)=\diag(1+z^{r_1},\ldots,1+z^{r_4}),\qquad
 A_\pm(z)=N_0(z)-N_\pm(z).
\]
The entries of $N_\pm$ have nonnegative integer coefficients. In entry $(i,j)$,
every exponent $p$ satisfies $0<p<r_i$, and the two signs have disjoint
support. Require the symplectic identity
\begin{equation}\label{eq:symplectic}
 A_+(z)A_-(z^{-1})^{\mathsf T}
 =A_-(z)A_+(z^{-1})^{\mathsf T}.
\end{equation}
The datum is indecomposable if no simultaneous index permutation makes both
matrices block diagonal with two nonempty blocks. Finite type means periodicity
of the associated universal $Y$-system, as in \cite[Definition 1.4]{Mizuno}.
\end{definition}

\begin{theorem}[Computer-assisted classification]\label{thm:classification}
An indecomposable datum in Definition~\ref{def:datum} is of finite type if and
only if, after an index permutation and possibly exchanging the signs,
\begin{equation}\label{eq:family}
 A_\pm(z)=E(z)B_\pm^{(k)}(z^\lambda)E(z)^{-1},\qquad
 E(z)=\diag(z^{s_1},z^{s_2},z^{s_3},1),
\end{equation}
for exactly one of the thirty-seven representatives in
Section~\ref{sec:catalogue}. Here $\lambda\in\Q_{>0}$ and
$s_1,s_2,s_3\in\Q$. The parameters are precisely those for which all transformed
exponents are integral and the support inequalities hold.
\end{theorem}

Thus a representative term $z^p$ in entry $(i,j)$ becomes
$z^{\lambda p+s_i-s_j}$, and $r_i=\lambda r_i^{(k)}$, with $s_4=0$.
This specifies all admissible polynomial lifts at arbitrarily large delays.
The thirty-seven representatives also belong to distinct classes under change
of slices in \cite[Section 3.1]{Mizuno}; see Section~\ref{sec:slices}.

\begin{corollary}\label{cor:positivity}
For the rank-four data of Definition~\ref{def:datum}, finite type is equivalent
to the existence of a row vector $v>0$ with
$vA_+(1)>0$ and $vA_-(1)>0$.
\end{corollary}

For decomposable data, apply the theorem blockwise. The remaining rank
partitions are $3+1$, $2+2$, $2+1+1$ and $1+1+1+1$, using the rank-three,
rank-two and rank-one classifications. Each component may have its own time
scale; a common period is a common multiple of the component periods.
The statements concern $D=I_4$ and diagonal $N_0$, rather than the full
definition allowing other symmetrizers and leading permutations.

\section{A finite reduction without a delay cutoff}\label{sec:reduction}
Write $P=N_+(1)$ and $M=N_-(1)$, so that $A_\pm(1)=2I_4-P,2I_4-M$.
The necessary simultaneous positivity theorem is
\cite[Lemma 3.3]{Mizuno}, citing the arbitrary-rank theorem
\cite[Theorem 5.5]{TData}. Hence a finite-type datum has a row vector $v>0$ with
\begin{equation}\label{eq:positive}
 v(2I_4-P)>0,\qquad v(2I_4-M)>0.
\end{equation}
In particular $P_{ii},M_{ii}\in\{0,1\}$, and both $2I_4-P$ and $2I_4-M$
have positive principal minors. Indeed, scaling by $v$ makes their transposes
strictly diagonally dominant Z-matrices. Every principal submatrix retains
this property and has positive determinant.

\begin{lemma}\label{lem:strong}
For an indecomposable datum, the graph with arrows $i\to j$ when $i\ne j$ and
$P_{ij}+M_{ij}>0$ is strongly connected.
\end{lemma}
\begin{proof}
If the graph is reducible, order its strongly connected components so that
$A_+$ and $A_-$ are block lower triangular. The matrices have constant term
$I_4$. Thus $K=A_+^{-1}A_-$ exists over $\Q(z)$ and is regular at zero.
Equation~\eqref{eq:symplectic} gives $K(z)=K(z^{-1})^{\mathsf T}$, so its
block lower triangular form is block diagonal.

Suppose that some entry between distinct blocks of $N_+$ or $N_-$ is nonzero,
and let $h>0$ be the smallest degree of such an entry. In
$K=I_4+A_+^{-1}(N_+-N_-)$, the coefficient of degree $h$ between blocks is
the corresponding coefficient of $N_+-N_-$. Every other contribution has a
crossing factor of degree at least $h$ and a further factor of positive
degree. Disjoint support makes one degree-$h$ crossing coefficient nonzero,
contradicting block diagonality of $K$.
\end{proof}

\begin{lemma}[Coefficient bound]\label{lem:bound}
Every finite-type indecomposable datum satisfies $P_{ij},M_{ij}\le15$ for
$i\ne j$.
\end{lemma}
\begin{proof}
Let $w_{ij}=\max(P_{ij},M_{ij})$. Equation~\eqref{eq:positive} gives
$w_{ij}v_i<2v_j$. Every arrow belongs to a simple directed cycle of length
$\ell\le4$, by Lemma~\ref{lem:strong}. Multiplication along the cycle gives
$\prod w_{ij}<2^\ell$. All its weights are positive integers, so each is
at most $2^\ell-1\le15$.
\end{proof}

\begin{lemma}[Parity]\label{lem:parity}
Each diagonal entry of $(2I_4-P)(2I_4-M)^{\mathsf T}$ is even.
\end{lemma}
\begin{proof}
The diagonal Laurent polynomial $(A_+A_-^*)_{ii}$ is reciprocal, where
$F^*(z)=F(z^{-1})^{\mathsf T}$. Its constant coefficient is $2$: the leading
diagonal factors contribute $2$, the mixed leading/support terms contribute
zero, and disjoint support gives
$[z^0]\sum_j(N_+)_{ij}(z)(N_-)_{ij}(z^{-1})=0$.
Every nonconstant coefficient of a reciprocal Laurent polynomial occurs in
an equal pair. Evaluation at $1$ is therefore even.
\end{proof}

\subsection*{How the constants are enumerated}
Relabel a vector in \eqref{eq:positive} so that $v_1\ge v_2\ge v_3\ge v_4$.
For either $C=P$ or $C=M$, its $j$th column then satisfies
\[
 \sum_{i<j}C_{ij}\le 1-C_{jj}.
\]
Thus, above the diagonal, column $j$ is either zero or has a single $1$;
the latter requires $C_{jj}=0$. There are
$\prod_{j=1}^4(j+1)=120$ choices of the diagonal and strict upper triangle
of each matrix. Exchanging the signs leaves $120\cdot121/2=7,260$ unordered
pairs of these choices to process.

The program completes the rows in order. When row $j$ is reached, it
enumerates the unknown entries of $P$ in $\{0,\ldots,15\}$ and uses
\[
 (2I_4-P)(2I_4-M)^{\mathsf T}
 =(2I_4-M)(2I_4-P)^{\mathsf T}
\]
to solve uniquely for the unknown entries of that row of $M$. The coefficient
matrix is the already completed leading principal block of $2I_4-P$,
whose determinant is positive. Adjugates, divisibility checks, cycle bounds
and principal minors are computed with integers. No solution can be lost
by solving this nonsingular linear system.

Completed pairs are tested for strong connectivity and ordered simultaneous
positivity, and then identified under all $24$ index permutations and sign
exchange. For positivity, the C++ enumerator constructs extreme rays of
\[
 \{v:\ v\ge0,\ v(2I_4-P)\ge0,\ v(2I_4-M)\ge0,
             \ v_i\ge v_{i+1}\}.
\]
Each extreme ray is obtained from three independent active constraint
normals. The cone is pointed; the sum of its rays is strict in the required
inequalities exactly when a suitable $v$ exists. A separate rational linear
solver supplies and verifies an integer positive vector for every retained
pair, without imposing the ordering on its canonical representative.

\begin{computation}\label{comp:constants}
The enumeration gives $4,865$ distinct constant pairs. The independent
verifier checks their symplectic identities, all principal minors, strong
connectivity, canonical forms and explicit simultaneous positive vectors.
Lemma~\ref{lem:parity} excludes $4,405$ pairs, leaving $460$.
\end{computation}

As a control on completeness of the new enumeration method, its rank-three
version reproduces exactly the entire set of $180$ constant pairs in the
rank-three calculation, not merely its cardinality.

\section{The chain obstruction}\label{sec:chain}
The following elementary exclusion is used before the remaining feasibility
checks. It also gives a direct proof for the largest surviving coefficient
pattern in the constant enumeration.

\begin{lemma}\label{lem:chain}
No datum of Definition~\ref{def:datum} has $P,M$ of the form
\[
 P=\begin{pmatrix}0&0&0&0\\0&0&0&0\\0&0&0&0\\u&v&w&0\end{pmatrix},\qquad
 M=\begin{pmatrix}0&1&0&1\\1&0&1&0\\0&1&0&0\\0&0&t&0\end{pmatrix},
 \qquad u,v,w,t\in\Z_{\ge0}.
\]
\end{lemma}
\begin{proof}
The $(1,2)$ and $(2,3)$ symplectic equations first imply
$r_1=r_2=r_3=R$ and
\[
 (N_-)_{12}=z^a,\quad (N_-)_{21}=z^{R-a},\quad
 (N_-)_{23}=z^c,\quad (N_-)_{32}=z^{R-c}.
\]
For example, the first equation is equality of two translated binomials;
their lowest and highest exponents give the asserted equalities.
Write $(N_-)_{14}=z^b$, and define the Laurent polynomials
\[
 \begin{aligned}
 q&=z^R,&F&=1+q,&H&=z^b(1+z^{-r_4}),\\
 X&=(N_+)_{41}(z^{-1}),&
 Y&=z^a(N_+)_{42}(z^{-1}),\\
 Z&=z^{a+c}(N_+)_{43}(z^{-1}),&
 W&=z^{a+c}(N_-)_{43}(z^{-1}).
 \end{aligned}
\]
The $(1,4)$, $(2,4)$ and $(3,4)$ equations give, respectively,
\begin{equation}\label{eq:chain}
 Y=FX+H,\qquad FY=qX+Z,\qquad F(Z-W)=qY.
\end{equation}
Consequently $Z=(1+q+q^2)X+FH$. The polynomials $X,H$ have nonnegative
coefficients and $H\ne0$. The first two expressions therefore show that
$Y$ and $Z$ have the same lowest exponent $d$. Since $R>0$, the lowest
exponent of $qY$ is $d+R>d$. The last equation in \eqref{eq:chain} forces
$W$ to have lowest exponent $d$, with the same nonzero coefficient there
as $Z$. This gives a common exponent in $(N_+)_{43}$ and $(N_-)_{43}$,
contradicting disjoint support.
\end{proof}

The lemma applies to constant IDs $22,35,44$. ID $35$ has already been
removed by parity. Thus $458$ pairs remain for the unbounded polynomial
calculation.

\section{All polynomial lifts and their periods}\label{sec:lifts}
For each constant pair, introduce four delay variables and one exponent
variable for every unit of coefficient mass. Equal exponents of the same
sign are allowed, so coefficients greater than one are included. Exponents
within an entry are ordered, opposite signs have distinct exponents, and
every exponent is strictly between zero and its row delay.

Expand $A_+A_-^*-A_-A_+^*$. In each entry, equality to zero is equality of
two finite multisets of integral linear forms in these variables. The main
encoding compares multiplicities at every positive-side exponent; equality
of the total sizes follows from the constant-matrix identity. A second
encoding sorts the two lists through exact comparison networks and equates
them. Both describe the full polynomial identity, including coincident
monomials. There is no upper bound on delays or exponents.

Most queries use integer variables. A homogeneous query can also be solved
over the reals with $\sum r_i=1$. Each Boolean branch is a rational linear
system with strict inequalities. A nonempty such system has a rational
point; clearing denominators yields an integer solution of the original
homogeneous query. Conversely any integer solution can be normalized.
This justifies the rational arithmetic used for the exclusion of ID $1117$.

\begin{computation}\label{comp:lifts}
Of the $458$ remaining constant pairs, $421$ have no polynomial lift and
$37$ have a lift. For each survivor, all admissible lifts lie in a single
four-dimensional rational linear subspace. Its generators are the
representative's exponent vector and the three relative species shifts in
\eqref{eq:family}. There are no unresolved solver answers.
\end{computation}

To obtain a subspace, match equal exponents on the two sides of each Laurent
identity at a satisfying point. The matching imposes homogeneous linear
equations that imply the entire identity. The verifier reduces the exponent
forms modulo the resulting row-reduced matrix and checks equality of the
two multisets symbolically. It also checks the dimension and the four
generators, then asks whether any admissible lift lies outside that subspace.
Every such coverage query is UNSAT. This proves necessity in
Theorem~\ref{thm:classification} for every delay.

\subsection*{Sufficiency}
For each representative, reconstruct the skew-symmetric exchange matrix
on $\{(i,p):1\le i\le4,\ 0\le p<r_i\}$ using
\cite[equation (2.2)]{Mizuno}. A time step mutates the four vertices $(i,0)$
and relabels $(i,p)\mapsto(i,p-1\bmod r_i)$. The verifier checks their
commutativity, the return of the exchange matrix and inverse-step consistency.
Starting from the identity principal C-matrix, integer mutation gives a
negative permutation matrix in both directions for all $37$ representatives.
Proposition 2.7 of \cite{Mizuno} then proves finite type in every semifield.
The lengths $h_+,h_-$ and the first positive return $\Omega$ of the labelled
C-matrix are listed in Section~\ref{sec:periods}. These are exact seed
certificates, rather than numerical tests of a few initial values.

Integral species shifts reindex solutions, and substituting $z^m$ splits
time into $m$ independent residue classes. Both operations preserve finite
type in both directions. For rational parameters in \eqref{eq:family}, clear
their denominators by an integer $L$. The transformed datum at $z^L$ is an
integral species shift of the representative at $z^{L\lambda}$. The integral
argument therefore proves sufficiency throughout each family. Combined
with \eqref{eq:positive} and Computation~\ref{comp:lifts}, this also proves
Corollary~\ref{cor:positivity}.

\section{Equivalences and the degree-two check}\label{sec:slices}
Specialization at $z=1$ is invariant under rescaling and species shifts.
The $37$ canonical constant pairs are distinct under index permutations
and sign exchange. Thus the explicit families in \eqref{eq:family} do not
overlap.

For change of slices, integral dilations repeat a slice sequence, while
species shifts move mutation events and the initial cut. The admissible
shift domain is convex. Along a straight path of shifts, the event order
changes only at coincident times; after clearing denominators these are
simultaneous mutations in a valid datum and they commute. Moving the initial
cut rotates the sequence. Rational rescalings and admissible shifts thus
give changes of slices.

Distinctness under this larger description is checked directly on connected
quiver slices. For each representative, follow one connected component
until the time-step permutation returns it. Its mutation word has four
events. The canonical encoding records the exchange matrix and the cycles
of the terminal permutation rooted at those events. The program exhausts
cyclic rotations and adjacent commuting mutations, including the negated
exchange matrix for sign exchange. All $37$ canonical minima are distinct.
Full encodings are stored, together with hashes; distinctness is checked
on the encodings themselves.

\begin{computation}[All delays equal to two]\label{comp:degree2}
The degree-two subcase has nine classes. Enumeration of all $3^{10}=59,049$
symmetric signed matrices gives $190$ connected labelled pairs satisfying
commutation and positive principal minors, and nine orbits. Each has an
exact two-sided reddening certificate. They occur in the full catalogue
as classes $1,2,3,28,29,30,31,34,35$.
\end{computation}

For this subcase $N_\pm(z)=zC_\pm$. Comparing Laurent coefficients forces
$C_+-C_-$ to be symmetric; disjoint support then makes each $C_\pm$
symmetric. The remaining identity is $C_+C_-=C_-C_+$. Positivity bounds
every entry by one, so the signed symmetric matrix has ten independent
entries in $\{-1,0,1\}$. This proves that the degree-two check is exhaustive.

\clearpage
\section{Periods and slice sizes}\label{sec:periods}
The class number is used throughout the catalogue; ID is the index in the
exhaustive constant enumeration. The period $\Omega$ is measured in the
displayed time coordinate. A slice is one connected component of the
reconstructed exchange matrix.
\begin{center}\small\renewcommand{\arraystretch}{1.1}
\begin{tabular}{@{}crcrrrr@{}}
\toprule
Class & ID & Delays & $h_+$ & $h_-$ & $\Omega$ & Slice size\\
\midrule
1 & 1 & $(2,2,2,2)$ & 2 & 6 & 8 & 4\\
2 & 2 & $(2,2,2,2)$ & 2 & 5 & 14 & 4\\
3 & 3 & $(2,2,2,2)$ & 2 & 9 & 11 & 8\\
4 & 4 & $(3,3,3,2)$ & 3 & 11 & 14 & 11\\
5 & 5 & $(6,6,2,6)$ & 6 & 20 & 26 & 10\\
6 & 36 & $(2,2,2,14)$ & 14 & 30 & 44 & 10\\
7 & 46 & $(5,5,2,2)$ & 5 & 12 & 17 & 14\\
8 & 47 & $(8,8,2,6)$ & 8 & 22 & 30 & 12\\
9 & 48 & $(5,5,2,3)$ & 5 & 13 & 18 & 15\\
10 & 56 & $(10,2,2,8)$ & 10 & 18 & 28 & 11\\
11 & 57 & $(7,2,2,5)$ & 7 & 12 & 19 & 16\\
12 & 58 & $(10,2,2,8)$ & 10 & 22 & 32 & 11\\
13 & 65 & $(2,2,2,4)$ & 6 & 4 & 20 & 5\\
14 & 79 & $(2,2,2,10)$ & 12 & 10 & 22 & 8\\
15 & 82 & $(2,2,2,7)$ & 9 & 7 & 16 & 13\\
16 & 122 & $(3,3,2,2)$ & 3 & 15 & 18 & 10\\
17 & 136 & $(2,2,8,8)$ & 12 & 30 & 42 & 10\\
18 & 207 & $(14,2,6,8)$ & 14 & 22 & 36 & 15\\
19 & 413 & $(2,2,16,12)$ & 42 & 24 & 66 & 16\\
20 & 432 & $(8,8,2,2)$ & 8 & 18 & 26 & 10\\
21 & 436 & $(14,6,2,6)$ & 14 & 20 & 34 & 14\\
22 & 437 & $(8,3,2,3)$ & 8 & 11 & 19 & 16\\
23 & 453 & $(2,12,2,12)$ & 26 & 18 & 44 & 14\\
24 & 456 & $(7,2,3,2)$ & 7 & 15 & 22 & 14\\
25 & 469 & $(8,3,2,5)$ & 8 & 13 & 21 & 18\\
26 & 949 & $(2,12,10,12)$ & 34 & 18 & 52 & 18\\
27 & 1262 & $(3,2,2,2)$ & 3 & 9 & 12 & 9\\
28 & 1493 & $(2,2,2,2)$ & 6 & 3 & 18 & 8\\
29 & 1494 & $(2,2,2,2)$ & 6 & 3 & 9 & 8\\
30 & 1834 & $(2,2,2,2)$ & 3 & 3 & 12 & 4\\
31 & 1836 & $(2,2,2,2)$ & 3 & 5 & 16 & 8\\
32 & 1844 & $(2,2,2,6)$ & 8 & 8 & 16 & 6\\
33 & 1848 & $(2,2,2,10)$ & 18 & 18 & 36 & 8\\
34 & 1852 & $(2,2,2,2)$ & 5 & 5 & 20 & 8\\
35 & 1853 & $(2,2,2,2)$ & 5 & 3 & 16 & 8\\
36 & 1902 & $(2,2,2,3)$ & 5 & 5 & 10 & 9\\
37 & 2614 & $(2,2,6,6)$ & 12 & 12 & 24 & 8\\
\bottomrule
\end{tabular}

\end{center}
\vfill

\clearpage
\section{The thirty-seven representatives}\label{sec:catalogue}
Each class displays $N_+^{(k)}(z),N_-^{(k)}(z)$ and the full constant matrices
\[
 A_\pm(1)=B_\pm^{(k)}(1)=2I_4-N_\pm^{(k)}(1).
\]
These constants are unchanged throughout the corresponding family;
relabelling simultaneously permutes rows and columns, and sign exchange
interchanges the pair. All entries below are generated from the checked
integer witnesses.
\subsection*{Classes 1--2}
\noindent $B_\pm^{(k)}(z)=\operatorname{diag}(1+z^{r_1},\ldots,1+z^{r_4})-N_\pm^{(k)}(z)$.\par\medskip
\noindent\begin{minipage}{\linewidth}
\noindent\textbf{Class 1}\quad\small (constant ID 1)\hfill $\symbf{r}=(2,2,2,2)$
\par\smallskip\noindent\textcolor{rulegray}{\rule{\linewidth}{0.3pt}}
\begin{minipage}[c]{0.495\linewidth}\centering\fontsize{9.7}{12}\selectfont
\setlength{\abovedisplayskip}{4pt}\setlength{\belowdisplayskip}{4pt}
\setlength{\abovedisplayshortskip}{4pt}\setlength{\belowdisplayshortskip}{4pt}
\setlength{\arraycolsep}{3pt}\renewcommand{\arraystretch}{1.08}
\[N_{+}^{(1)}(z)=\begin{pmatrix}0 & 0 & 0 & 0\\0 & 0 & 0 & 0\\0 & 0 & 0 & 0\\0 & 0 & 0 & 0\end{pmatrix}\]
\end{minipage}%
\begin{minipage}[c]{0.495\linewidth}\centering\fontsize{9.7}{12}\selectfont
\setlength{\abovedisplayskip}{4pt}\setlength{\belowdisplayskip}{4pt}
\setlength{\abovedisplayshortskip}{4pt}\setlength{\belowdisplayshortskip}{4pt}
\setlength{\arraycolsep}{3pt}\renewcommand{\arraystretch}{1.08}
\[N_{-}^{(1)}(z)=\begin{pmatrix}0 & 0 & 0 & z\\0 & 0 & 0 & z\\0 & 0 & 0 & z\\z & z & z & 0\end{pmatrix}\]
\end{minipage}%
\par\noindent
\begin{minipage}[c]{0.495\linewidth}\centering\fontsize{9.7}{12}\selectfont
\setlength{\abovedisplayskip}{4pt}\setlength{\belowdisplayskip}{4pt}
\setlength{\abovedisplayshortskip}{4pt}\setlength{\belowdisplayshortskip}{4pt}
\setlength{\arraycolsep}{3pt}\renewcommand{\arraystretch}{1.08}
\[A_{+}(1)=\begin{pmatrix}2 & 0 & 0 & 0\\0 & 2 & 0 & 0\\0 & 0 & 2 & 0\\0 & 0 & 0 & 2\end{pmatrix}\]
\end{minipage}%
\begin{minipage}[c]{0.495\linewidth}\centering\fontsize{9.7}{12}\selectfont
\setlength{\abovedisplayskip}{4pt}\setlength{\belowdisplayskip}{4pt}
\setlength{\abovedisplayshortskip}{4pt}\setlength{\belowdisplayshortskip}{4pt}
\setlength{\arraycolsep}{3pt}\renewcommand{\arraystretch}{1.08}
\[A_{-}(1)=\begin{pmatrix}2 & 0 & 0 & -1\\0 & 2 & 0 & -1\\0 & 0 & 2 & -1\\-1 & -1 & -1 & 2\end{pmatrix}\]
\end{minipage}%
\par\smallskip\noindent\small $h_+=2,\quad h_-=6,\quad \Omega=8.$\quad Vertices per slice: 4.
\end{minipage}\par\vfill
\noindent\begin{minipage}{\linewidth}
\noindent\textbf{Class 2}\quad\small (constant ID 2)\hfill $\symbf{r}=(2,2,2,2)$
\par\smallskip\noindent\textcolor{rulegray}{\rule{\linewidth}{0.3pt}}
\begin{minipage}[c]{0.495\linewidth}\centering\fontsize{9.7}{12}\selectfont
\setlength{\abovedisplayskip}{4pt}\setlength{\belowdisplayskip}{4pt}
\setlength{\abovedisplayshortskip}{4pt}\setlength{\belowdisplayshortskip}{4pt}
\setlength{\arraycolsep}{3pt}\renewcommand{\arraystretch}{1.08}
\[N_{+}^{(2)}(z)=\begin{pmatrix}0 & 0 & 0 & 0\\0 & 0 & 0 & 0\\0 & 0 & 0 & 0\\0 & 0 & 0 & 0\end{pmatrix}\]
\end{minipage}%
\begin{minipage}[c]{0.495\linewidth}\centering\fontsize{9.7}{12}\selectfont
\setlength{\abovedisplayskip}{4pt}\setlength{\belowdisplayskip}{4pt}
\setlength{\abovedisplayshortskip}{4pt}\setlength{\belowdisplayshortskip}{4pt}
\setlength{\arraycolsep}{3pt}\renewcommand{\arraystretch}{1.08}
\[N_{-}^{(2)}(z)=\begin{pmatrix}0 & 0 & 0 & z\\0 & 0 & z & 0\\0 & z & 0 & z\\z & 0 & z & 0\end{pmatrix}\]
\end{minipage}%
\par\noindent
\begin{minipage}[c]{0.495\linewidth}\centering\fontsize{9.7}{12}\selectfont
\setlength{\abovedisplayskip}{4pt}\setlength{\belowdisplayskip}{4pt}
\setlength{\abovedisplayshortskip}{4pt}\setlength{\belowdisplayshortskip}{4pt}
\setlength{\arraycolsep}{3pt}\renewcommand{\arraystretch}{1.08}
\[A_{+}(1)=\begin{pmatrix}2 & 0 & 0 & 0\\0 & 2 & 0 & 0\\0 & 0 & 2 & 0\\0 & 0 & 0 & 2\end{pmatrix}\]
\end{minipage}%
\begin{minipage}[c]{0.495\linewidth}\centering\fontsize{9.7}{12}\selectfont
\setlength{\abovedisplayskip}{4pt}\setlength{\belowdisplayskip}{4pt}
\setlength{\abovedisplayshortskip}{4pt}\setlength{\belowdisplayshortskip}{4pt}
\setlength{\arraycolsep}{3pt}\renewcommand{\arraystretch}{1.08}
\[A_{-}(1)=\begin{pmatrix}2 & 0 & 0 & -1\\0 & 2 & -1 & 0\\0 & -1 & 2 & -1\\-1 & 0 & -1 & 2\end{pmatrix}\]
\end{minipage}%
\par\smallskip\noindent\small $h_+=2,\quad h_-=5,\quad \Omega=14.$\quad Vertices per slice: 4.
\end{minipage}\par\vfill
\clearpage
\subsection*{Classes 3--5}
\noindent $B_\pm^{(k)}(z)=\operatorname{diag}(1+z^{r_1},\ldots,1+z^{r_4})-N_\pm^{(k)}(z)$.\par\medskip
\noindent\begin{minipage}{\linewidth}
\noindent\textbf{Class 3}\quad\small (constant ID 3)\hfill $\symbf{r}=(2,2,2,2)$
\par\smallskip\noindent\textcolor{rulegray}{\rule{\linewidth}{0.3pt}}
\begin{minipage}[c]{0.495\linewidth}\centering\fontsize{9.7}{12}\selectfont
\setlength{\abovedisplayskip}{4pt}\setlength{\belowdisplayskip}{4pt}
\setlength{\abovedisplayshortskip}{4pt}\setlength{\belowdisplayshortskip}{4pt}
\setlength{\arraycolsep}{3pt}\renewcommand{\arraystretch}{1.08}
\[N_{+}^{(3)}(z)=\begin{pmatrix}0 & 0 & 0 & 0\\0 & 0 & 0 & 0\\0 & 0 & 0 & 0\\0 & 0 & 0 & 0\end{pmatrix}\]
\end{minipage}%
\begin{minipage}[c]{0.495\linewidth}\centering\fontsize{9.7}{12}\selectfont
\setlength{\abovedisplayskip}{4pt}\setlength{\belowdisplayskip}{4pt}
\setlength{\abovedisplayshortskip}{4pt}\setlength{\belowdisplayshortskip}{4pt}
\setlength{\arraycolsep}{3pt}\renewcommand{\arraystretch}{1.08}
\[N_{-}^{(3)}(z)=\begin{pmatrix}0 & 0 & 0 & z\\0 & 0 & z & z\\0 & z & z & 0\\z & z & 0 & 0\end{pmatrix}\]
\end{minipage}%
\par\noindent
\begin{minipage}[c]{0.495\linewidth}\centering\fontsize{9.7}{12}\selectfont
\setlength{\abovedisplayskip}{4pt}\setlength{\belowdisplayskip}{4pt}
\setlength{\abovedisplayshortskip}{4pt}\setlength{\belowdisplayshortskip}{4pt}
\setlength{\arraycolsep}{3pt}\renewcommand{\arraystretch}{1.08}
\[A_{+}(1)=\begin{pmatrix}2 & 0 & 0 & 0\\0 & 2 & 0 & 0\\0 & 0 & 2 & 0\\0 & 0 & 0 & 2\end{pmatrix}\]
\end{minipage}%
\begin{minipage}[c]{0.495\linewidth}\centering\fontsize{9.7}{12}\selectfont
\setlength{\abovedisplayskip}{4pt}\setlength{\belowdisplayskip}{4pt}
\setlength{\abovedisplayshortskip}{4pt}\setlength{\belowdisplayshortskip}{4pt}
\setlength{\arraycolsep}{3pt}\renewcommand{\arraystretch}{1.08}
\[A_{-}(1)=\begin{pmatrix}2 & 0 & 0 & -1\\0 & 2 & -1 & -1\\0 & -1 & 1 & 0\\-1 & -1 & 0 & 2\end{pmatrix}\]
\end{minipage}%
\par\smallskip\noindent\small $h_+=2,\quad h_-=9,\quad \Omega=11.$\quad Vertices per slice: 8.
\end{minipage}\par\vfill
\noindent\begin{minipage}{\linewidth}
\noindent\textbf{Class 4}\quad\small (constant ID 4)\hfill $\symbf{r}=(3,3,3,2)$
\par\smallskip\noindent\textcolor{rulegray}{\rule{\linewidth}{0.3pt}}
\begin{minipage}[c]{0.495\linewidth}\centering\fontsize{9.7}{12}\selectfont
\setlength{\abovedisplayskip}{4pt}\setlength{\belowdisplayskip}{4pt}
\setlength{\abovedisplayshortskip}{4pt}\setlength{\belowdisplayshortskip}{4pt}
\setlength{\arraycolsep}{3pt}\renewcommand{\arraystretch}{1.08}
\[N_{+}^{(4)}(z)=\begin{pmatrix}0 & 0 & 0 & 0\\0 & 0 & 0 & 0\\0 & 0 & 0 & 0\\0 & 0 & 0 & z\end{pmatrix}\]
\end{minipage}%
\begin{minipage}[c]{0.495\linewidth}\centering\fontsize{9.7}{12}\selectfont
\setlength{\abovedisplayskip}{4pt}\setlength{\belowdisplayskip}{4pt}
\setlength{\abovedisplayshortskip}{4pt}\setlength{\belowdisplayshortskip}{4pt}
\setlength{\arraycolsep}{3pt}\renewcommand{\arraystretch}{1.08}
\[N_{-}^{(4)}(z)=\begin{pmatrix}0 & 0 & z & 0\\0 & 0 & z^{2} & z^{2} + z\\z^{2} & z & 0 & 0\\0 & z & 0 & 0\end{pmatrix}\]
\end{minipage}%
\par\noindent
\begin{minipage}[c]{0.495\linewidth}\centering\fontsize{9.7}{12}\selectfont
\setlength{\abovedisplayskip}{4pt}\setlength{\belowdisplayskip}{4pt}
\setlength{\abovedisplayshortskip}{4pt}\setlength{\belowdisplayshortskip}{4pt}
\setlength{\arraycolsep}{3pt}\renewcommand{\arraystretch}{1.08}
\[A_{+}(1)=\begin{pmatrix}2 & 0 & 0 & 0\\0 & 2 & 0 & 0\\0 & 0 & 2 & 0\\0 & 0 & 0 & 1\end{pmatrix}\]
\end{minipage}%
\begin{minipage}[c]{0.495\linewidth}\centering\fontsize{9.7}{12}\selectfont
\setlength{\abovedisplayskip}{4pt}\setlength{\belowdisplayskip}{4pt}
\setlength{\abovedisplayshortskip}{4pt}\setlength{\belowdisplayshortskip}{4pt}
\setlength{\arraycolsep}{3pt}\renewcommand{\arraystretch}{1.08}
\[A_{-}(1)=\begin{pmatrix}2 & 0 & -1 & 0\\0 & 2 & -1 & -2\\-1 & -1 & 2 & 0\\0 & -1 & 0 & 2\end{pmatrix}\]
\end{minipage}%
\par\smallskip\noindent\small $h_+=3,\quad h_-=11,\quad \Omega=14.$\quad Vertices per slice: 11.
\end{minipage}\par\vfill
\noindent\begin{minipage}{\linewidth}
\noindent\textbf{Class 5}\quad\small (constant ID 5)\hfill $\symbf{r}=(6,6,2,6)$
\par\smallskip\noindent\textcolor{rulegray}{\rule{\linewidth}{0.3pt}}
\begin{minipage}[c]{0.495\linewidth}\centering\fontsize{9.7}{12}\selectfont
\setlength{\abovedisplayskip}{4pt}\setlength{\belowdisplayskip}{4pt}
\setlength{\abovedisplayshortskip}{4pt}\setlength{\belowdisplayshortskip}{4pt}
\setlength{\arraycolsep}{3pt}\renewcommand{\arraystretch}{1.08}
\[N_{+}^{(5)}(z)=\begin{pmatrix}0 & 0 & 0 & 0\\0 & 0 & 0 & 0\\0 & 0 & 0 & 0\\0 & 0 & z^{3} & 0\end{pmatrix}\]
\end{minipage}%
\begin{minipage}[c]{0.495\linewidth}\centering\fontsize{9.7}{12}\selectfont
\setlength{\abovedisplayskip}{4pt}\setlength{\belowdisplayskip}{4pt}
\setlength{\abovedisplayshortskip}{4pt}\setlength{\belowdisplayshortskip}{4pt}
\setlength{\arraycolsep}{3pt}\renewcommand{\arraystretch}{1.08}
\[N_{-}^{(5)}(z)=\begin{pmatrix}0 & z & 0 & 0\\z^{5} & 0 & 0 & z\\0 & 0 & 0 & z\\0 & z^{5} & z^{5} + z & 0\end{pmatrix}\]
\end{minipage}%
\par\noindent
\begin{minipage}[c]{0.495\linewidth}\centering\fontsize{9.7}{12}\selectfont
\setlength{\abovedisplayskip}{4pt}\setlength{\belowdisplayskip}{4pt}
\setlength{\abovedisplayshortskip}{4pt}\setlength{\belowdisplayshortskip}{4pt}
\setlength{\arraycolsep}{3pt}\renewcommand{\arraystretch}{1.08}
\[A_{+}(1)=\begin{pmatrix}2 & 0 & 0 & 0\\0 & 2 & 0 & 0\\0 & 0 & 2 & 0\\0 & 0 & -1 & 2\end{pmatrix}\]
\end{minipage}%
\begin{minipage}[c]{0.495\linewidth}\centering\fontsize{9.7}{12}\selectfont
\setlength{\abovedisplayskip}{4pt}\setlength{\belowdisplayskip}{4pt}
\setlength{\abovedisplayshortskip}{4pt}\setlength{\belowdisplayshortskip}{4pt}
\setlength{\arraycolsep}{3pt}\renewcommand{\arraystretch}{1.08}
\[A_{-}(1)=\begin{pmatrix}2 & -1 & 0 & 0\\-1 & 2 & 0 & -1\\0 & 0 & 2 & -1\\0 & -1 & -2 & 2\end{pmatrix}\]
\end{minipage}%
\par\smallskip\noindent\small $h_+=6,\quad h_-=20,\quad \Omega=26.$\quad Vertices per slice: 10.
\end{minipage}\par\vfill
\clearpage
\subsection*{Classes 6--8}
\noindent $B_\pm^{(k)}(z)=\operatorname{diag}(1+z^{r_1},\ldots,1+z^{r_4})-N_\pm^{(k)}(z)$.\par\medskip
\noindent\begin{minipage}{\linewidth}
\noindent\textbf{Class 6}\quad\small (constant ID 36)\hfill $\symbf{r}=(2,2,2,14)$
\par\smallskip\noindent\textcolor{rulegray}{\rule{\linewidth}{0.3pt}}
\begin{minipage}[c]{0.495\linewidth}\centering\fontsize{9.7}{12}\selectfont
\setlength{\abovedisplayskip}{4pt}\setlength{\belowdisplayskip}{4pt}
\setlength{\abovedisplayshortskip}{4pt}\setlength{\belowdisplayshortskip}{4pt}
\setlength{\arraycolsep}{3pt}\renewcommand{\arraystretch}{1.08}
\[N_{+}^{(6)}(z)=\begin{pmatrix}0 & 0 & 0 & 0\\0 & 0 & 0 & 0\\0 & 0 & 0 & 0\\z^{9} + z^{5} & z^{11} + z^{3} & z^{10} + z^{4} & 0\end{pmatrix}\]
\end{minipage}%
\begin{minipage}[c]{0.495\linewidth}\centering\fontsize{9.7}{12}\selectfont
\setlength{\abovedisplayskip}{4pt}\setlength{\belowdisplayskip}{4pt}
\setlength{\abovedisplayshortskip}{4pt}\setlength{\belowdisplayshortskip}{4pt}
\setlength{\arraycolsep}{3pt}\renewcommand{\arraystretch}{1.08}
\[N_{-}^{(6)}(z)=\begin{pmatrix}0 & 0 & z & 0\\0 & 0 & z & z\\z & z & 0 & 0\\z^{7} & z^{13} + z & 0 & 0\end{pmatrix}\]
\end{minipage}%
\par\noindent
\begin{minipage}[c]{0.495\linewidth}\centering\fontsize{9.7}{12}\selectfont
\setlength{\abovedisplayskip}{4pt}\setlength{\belowdisplayskip}{4pt}
\setlength{\abovedisplayshortskip}{4pt}\setlength{\belowdisplayshortskip}{4pt}
\setlength{\arraycolsep}{3pt}\renewcommand{\arraystretch}{1.08}
\[A_{+}(1)=\begin{pmatrix}2 & 0 & 0 & 0\\0 & 2 & 0 & 0\\0 & 0 & 2 & 0\\-2 & -2 & -2 & 2\end{pmatrix}\]
\end{minipage}%
\begin{minipage}[c]{0.495\linewidth}\centering\fontsize{9.7}{12}\selectfont
\setlength{\abovedisplayskip}{4pt}\setlength{\belowdisplayskip}{4pt}
\setlength{\abovedisplayshortskip}{4pt}\setlength{\belowdisplayshortskip}{4pt}
\setlength{\arraycolsep}{3pt}\renewcommand{\arraystretch}{1.08}
\[A_{-}(1)=\begin{pmatrix}2 & 0 & -1 & 0\\0 & 2 & -1 & -1\\-1 & -1 & 2 & 0\\-1 & -2 & 0 & 2\end{pmatrix}\]
\end{minipage}%
\par\smallskip\noindent\small $h_+=14,\quad h_-=30,\quad \Omega=44.$\quad Vertices per slice: 10.
\end{minipage}\par\vfill
\noindent\begin{minipage}{\linewidth}
\noindent\textbf{Class 7}\quad\small (constant ID 46)\hfill $\symbf{r}=(5,5,2,2)$
\par\smallskip\noindent\textcolor{rulegray}{\rule{\linewidth}{0.3pt}}
\begin{minipage}[c]{0.495\linewidth}\centering\fontsize{9.7}{12}\selectfont
\setlength{\abovedisplayskip}{4pt}\setlength{\belowdisplayskip}{4pt}
\setlength{\abovedisplayshortskip}{4pt}\setlength{\belowdisplayshortskip}{4pt}
\setlength{\arraycolsep}{3pt}\renewcommand{\arraystretch}{1.08}
\[N_{+}^{(7)}(z)=\begin{pmatrix}0 & 0 & 0 & 0\\0 & 0 & 0 & 0\\0 & 0 & 0 & z\\0 & 0 & z & z\end{pmatrix}\]
\end{minipage}%
\begin{minipage}[c]{0.495\linewidth}\centering\fontsize{9.7}{12}\selectfont
\setlength{\abovedisplayskip}{4pt}\setlength{\belowdisplayskip}{4pt}
\setlength{\abovedisplayshortskip}{4pt}\setlength{\belowdisplayshortskip}{4pt}
\setlength{\arraycolsep}{3pt}\renewcommand{\arraystretch}{1.08}
\[N_{-}^{(7)}(z)=\begin{pmatrix}0 & z & 0 & 0\\z^{4} & 0 & z^{4} + z & z^{3} + z^{2}\\0 & z & 0 & 0\\0 & 0 & 0 & 0\end{pmatrix}\]
\end{minipage}%
\par\noindent
\begin{minipage}[c]{0.495\linewidth}\centering\fontsize{9.7}{12}\selectfont
\setlength{\abovedisplayskip}{4pt}\setlength{\belowdisplayskip}{4pt}
\setlength{\abovedisplayshortskip}{4pt}\setlength{\belowdisplayshortskip}{4pt}
\setlength{\arraycolsep}{3pt}\renewcommand{\arraystretch}{1.08}
\[A_{+}(1)=\begin{pmatrix}2 & 0 & 0 & 0\\0 & 2 & 0 & 0\\0 & 0 & 2 & -1\\0 & 0 & -1 & 1\end{pmatrix}\]
\end{minipage}%
\begin{minipage}[c]{0.495\linewidth}\centering\fontsize{9.7}{12}\selectfont
\setlength{\abovedisplayskip}{4pt}\setlength{\belowdisplayskip}{4pt}
\setlength{\abovedisplayshortskip}{4pt}\setlength{\belowdisplayshortskip}{4pt}
\setlength{\arraycolsep}{3pt}\renewcommand{\arraystretch}{1.08}
\[A_{-}(1)=\begin{pmatrix}2 & -1 & 0 & 0\\-1 & 2 & -2 & -2\\0 & -1 & 2 & 0\\0 & 0 & 0 & 2\end{pmatrix}\]
\end{minipage}%
\par\smallskip\noindent\small $h_+=5,\quad h_-=12,\quad \Omega=17.$\quad Vertices per slice: 14.
\end{minipage}\par\vfill
\noindent\begin{minipage}{\linewidth}
\noindent\textbf{Class 8}\quad\small (constant ID 47)\hfill $\symbf{r}=(8,8,2,6)$
\par\smallskip\noindent\textcolor{rulegray}{\rule{\linewidth}{0.3pt}}
\begin{minipage}[c]{0.495\linewidth}\centering\fontsize{9.7}{12}\selectfont
\setlength{\abovedisplayskip}{4pt}\setlength{\belowdisplayskip}{4pt}
\setlength{\abovedisplayshortskip}{4pt}\setlength{\belowdisplayshortskip}{4pt}
\setlength{\arraycolsep}{3pt}\renewcommand{\arraystretch}{1.08}
\[N_{+}^{(8)}(z)=\begin{pmatrix}0 & 0 & 0 & 0\\0 & 0 & 0 & 0\\0 & 0 & 0 & z\\0 & 0 & z^{5} + z & 0\end{pmatrix}\]
\end{minipage}%
\begin{minipage}[c]{0.495\linewidth}\centering\fontsize{9.7}{12}\selectfont
\setlength{\abovedisplayskip}{4pt}\setlength{\belowdisplayskip}{4pt}
\setlength{\abovedisplayshortskip}{4pt}\setlength{\belowdisplayshortskip}{4pt}
\setlength{\arraycolsep}{3pt}\renewcommand{\arraystretch}{1.08}
\[N_{-}^{(8)}(z)=\begin{pmatrix}0 & z^{7} & 0 & 0\\z & 0 & z^{2} & z^{3} + z\\0 & 0 & 0 & 0\\0 & z^{5} & z^{3} & 0\end{pmatrix}\]
\end{minipage}%
\par\noindent
\begin{minipage}[c]{0.495\linewidth}\centering\fontsize{9.7}{12}\selectfont
\setlength{\abovedisplayskip}{4pt}\setlength{\belowdisplayskip}{4pt}
\setlength{\abovedisplayshortskip}{4pt}\setlength{\belowdisplayshortskip}{4pt}
\setlength{\arraycolsep}{3pt}\renewcommand{\arraystretch}{1.08}
\[A_{+}(1)=\begin{pmatrix}2 & 0 & 0 & 0\\0 & 2 & 0 & 0\\0 & 0 & 2 & -1\\0 & 0 & -2 & 2\end{pmatrix}\]
\end{minipage}%
\begin{minipage}[c]{0.495\linewidth}\centering\fontsize{9.7}{12}\selectfont
\setlength{\abovedisplayskip}{4pt}\setlength{\belowdisplayskip}{4pt}
\setlength{\abovedisplayshortskip}{4pt}\setlength{\belowdisplayshortskip}{4pt}
\setlength{\arraycolsep}{3pt}\renewcommand{\arraystretch}{1.08}
\[A_{-}(1)=\begin{pmatrix}2 & -1 & 0 & 0\\-1 & 2 & -1 & -2\\0 & 0 & 2 & 0\\0 & -1 & -1 & 2\end{pmatrix}\]
\end{minipage}%
\par\smallskip\noindent\small $h_+=8,\quad h_-=22,\quad \Omega=30.$\quad Vertices per slice: 12.
\end{minipage}\par\vfill
\clearpage
\subsection*{Classes 9--11}
\noindent $B_\pm^{(k)}(z)=\operatorname{diag}(1+z^{r_1},\ldots,1+z^{r_4})-N_\pm^{(k)}(z)$.\par\medskip
\noindent\begin{minipage}{\linewidth}
\noindent\textbf{Class 9}\quad\small (constant ID 48)\hfill $\symbf{r}=(5,5,2,3)$
\par\smallskip\noindent\textcolor{rulegray}{\rule{\linewidth}{0.3pt}}
\begin{minipage}[c]{0.495\linewidth}\centering\fontsize{9.7}{12}\selectfont
\setlength{\abovedisplayskip}{4pt}\setlength{\belowdisplayskip}{4pt}
\setlength{\abovedisplayshortskip}{4pt}\setlength{\belowdisplayshortskip}{4pt}
\setlength{\arraycolsep}{3pt}\renewcommand{\arraystretch}{1.08}
\[N_{+}^{(9)}(z)=\begin{pmatrix}0 & 0 & 0 & 0\\0 & 0 & 0 & 0\\0 & 0 & 0 & z\\0 & 0 & z^{2} + z & 0\end{pmatrix}\]
\end{minipage}%
\begin{minipage}[c]{0.495\linewidth}\centering\fontsize{9.7}{12}\selectfont
\setlength{\abovedisplayskip}{4pt}\setlength{\belowdisplayskip}{4pt}
\setlength{\abovedisplayshortskip}{4pt}\setlength{\belowdisplayshortskip}{4pt}
\setlength{\arraycolsep}{3pt}\renewcommand{\arraystretch}{1.08}
\[N_{-}^{(9)}(z)=\begin{pmatrix}0 & z^{4} & 0 & 0\\z & 0 & z^{2} & z^{3} + z\\0 & 0 & z & 0\\0 & z^{2} & 0 & 0\end{pmatrix}\]
\end{minipage}%
\par\noindent
\begin{minipage}[c]{0.495\linewidth}\centering\fontsize{9.7}{12}\selectfont
\setlength{\abovedisplayskip}{4pt}\setlength{\belowdisplayskip}{4pt}
\setlength{\abovedisplayshortskip}{4pt}\setlength{\belowdisplayshortskip}{4pt}
\setlength{\arraycolsep}{3pt}\renewcommand{\arraystretch}{1.08}
\[A_{+}(1)=\begin{pmatrix}2 & 0 & 0 & 0\\0 & 2 & 0 & 0\\0 & 0 & 2 & -1\\0 & 0 & -2 & 2\end{pmatrix}\]
\end{minipage}%
\begin{minipage}[c]{0.495\linewidth}\centering\fontsize{9.7}{12}\selectfont
\setlength{\abovedisplayskip}{4pt}\setlength{\belowdisplayskip}{4pt}
\setlength{\abovedisplayshortskip}{4pt}\setlength{\belowdisplayshortskip}{4pt}
\setlength{\arraycolsep}{3pt}\renewcommand{\arraystretch}{1.08}
\[A_{-}(1)=\begin{pmatrix}2 & -1 & 0 & 0\\-1 & 2 & -1 & -2\\0 & 0 & 1 & 0\\0 & -1 & 0 & 2\end{pmatrix}\]
\end{minipage}%
\par\smallskip\noindent\small $h_+=5,\quad h_-=13,\quad \Omega=18.$\quad Vertices per slice: 15.
\end{minipage}\par\vfill
\noindent\begin{minipage}{\linewidth}
\noindent\textbf{Class 10}\quad\small (constant ID 56)\hfill $\symbf{r}=(10,2,2,8)$
\par\smallskip\noindent\textcolor{rulegray}{\rule{\linewidth}{0.3pt}}
\begin{minipage}[c]{0.495\linewidth}\centering\fontsize{9.7}{12}\selectfont
\setlength{\abovedisplayskip}{4pt}\setlength{\belowdisplayskip}{4pt}
\setlength{\abovedisplayshortskip}{4pt}\setlength{\belowdisplayshortskip}{4pt}
\setlength{\arraycolsep}{3pt}\renewcommand{\arraystretch}{1.08}
\[N_{+}^{(10)}(z)=\begin{pmatrix}0 & 0 & 0 & 0\\0 & 0 & 0 & 0\\0 & 0 & 0 & z\\0 & z^{6} + z^{2} & z^{7} + z & 0\end{pmatrix}\]
\end{minipage}%
\begin{minipage}[c]{0.495\linewidth}\centering\fontsize{9.7}{12}\selectfont
\setlength{\abovedisplayskip}{4pt}\setlength{\belowdisplayskip}{4pt}
\setlength{\abovedisplayshortskip}{4pt}\setlength{\belowdisplayshortskip}{4pt}
\setlength{\arraycolsep}{3pt}\renewcommand{\arraystretch}{1.08}
\[N_{-}^{(10)}(z)=\begin{pmatrix}0 & 0 & z^{2} & z^{3} + z\\0 & 0 & z & 0\\0 & z & 0 & 0\\z^{7} & z^{4} & 0 & 0\end{pmatrix}\]
\end{minipage}%
\par\noindent
\begin{minipage}[c]{0.495\linewidth}\centering\fontsize{9.7}{12}\selectfont
\setlength{\abovedisplayskip}{4pt}\setlength{\belowdisplayskip}{4pt}
\setlength{\abovedisplayshortskip}{4pt}\setlength{\belowdisplayshortskip}{4pt}
\setlength{\arraycolsep}{3pt}\renewcommand{\arraystretch}{1.08}
\[A_{+}(1)=\begin{pmatrix}2 & 0 & 0 & 0\\0 & 2 & 0 & 0\\0 & 0 & 2 & -1\\0 & -2 & -2 & 2\end{pmatrix}\]
\end{minipage}%
\begin{minipage}[c]{0.495\linewidth}\centering\fontsize{9.7}{12}\selectfont
\setlength{\abovedisplayskip}{4pt}\setlength{\belowdisplayskip}{4pt}
\setlength{\abovedisplayshortskip}{4pt}\setlength{\belowdisplayshortskip}{4pt}
\setlength{\arraycolsep}{3pt}\renewcommand{\arraystretch}{1.08}
\[A_{-}(1)=\begin{pmatrix}2 & 0 & -1 & -2\\0 & 2 & -1 & 0\\0 & -1 & 2 & 0\\-1 & -1 & 0 & 2\end{pmatrix}\]
\end{minipage}%
\par\smallskip\noindent\small $h_+=10,\quad h_-=18,\quad \Omega=28.$\quad Vertices per slice: 11.
\end{minipage}\par\vfill
\noindent\begin{minipage}{\linewidth}
\noindent\textbf{Class 11}\quad\small (constant ID 57)\hfill $\symbf{r}=(7,2,2,5)$
\par\smallskip\noindent\textcolor{rulegray}{\rule{\linewidth}{0.3pt}}
\begin{minipage}[c]{0.495\linewidth}\centering\fontsize{9.7}{12}\selectfont
\setlength{\abovedisplayskip}{4pt}\setlength{\belowdisplayskip}{4pt}
\setlength{\abovedisplayshortskip}{4pt}\setlength{\belowdisplayshortskip}{4pt}
\setlength{\arraycolsep}{3pt}\renewcommand{\arraystretch}{1.08}
\[N_{+}^{(11)}(z)=\begin{pmatrix}0 & 0 & 0 & 0\\0 & 0 & 0 & 0\\0 & 0 & 0 & z\\0 & z^{3} + z^{2} & z^{4} + z & 0\end{pmatrix}\]
\end{minipage}%
\begin{minipage}[c]{0.495\linewidth}\centering\fontsize{9.7}{12}\selectfont
\setlength{\abovedisplayskip}{4pt}\setlength{\belowdisplayskip}{4pt}
\setlength{\abovedisplayshortskip}{4pt}\setlength{\belowdisplayshortskip}{4pt}
\setlength{\arraycolsep}{3pt}\renewcommand{\arraystretch}{1.08}
\[N_{-}^{(11)}(z)=\begin{pmatrix}0 & 0 & z^{2} & z^{3} + z\\0 & z & z & 0\\0 & z & 0 & 0\\z^{4} & 0 & 0 & 0\end{pmatrix}\]
\end{minipage}%
\par\noindent
\begin{minipage}[c]{0.495\linewidth}\centering\fontsize{9.7}{12}\selectfont
\setlength{\abovedisplayskip}{4pt}\setlength{\belowdisplayskip}{4pt}
\setlength{\abovedisplayshortskip}{4pt}\setlength{\belowdisplayshortskip}{4pt}
\setlength{\arraycolsep}{3pt}\renewcommand{\arraystretch}{1.08}
\[A_{+}(1)=\begin{pmatrix}2 & 0 & 0 & 0\\0 & 2 & 0 & 0\\0 & 0 & 2 & -1\\0 & -2 & -2 & 2\end{pmatrix}\]
\end{minipage}%
\begin{minipage}[c]{0.495\linewidth}\centering\fontsize{9.7}{12}\selectfont
\setlength{\abovedisplayskip}{4pt}\setlength{\belowdisplayskip}{4pt}
\setlength{\abovedisplayshortskip}{4pt}\setlength{\belowdisplayshortskip}{4pt}
\setlength{\arraycolsep}{3pt}\renewcommand{\arraystretch}{1.08}
\[A_{-}(1)=\begin{pmatrix}2 & 0 & -1 & -2\\0 & 1 & -1 & 0\\0 & -1 & 2 & 0\\-1 & 0 & 0 & 2\end{pmatrix}\]
\end{minipage}%
\par\smallskip\noindent\small $h_+=7,\quad h_-=12,\quad \Omega=19.$\quad Vertices per slice: 16.
\end{minipage}\par\vfill
\clearpage
\subsection*{Classes 12--14}
\noindent $B_\pm^{(k)}(z)=\operatorname{diag}(1+z^{r_1},\ldots,1+z^{r_4})-N_\pm^{(k)}(z)$.\par\medskip
\noindent\begin{minipage}{\linewidth}
\noindent\textbf{Class 12}\quad\small (constant ID 58)\hfill $\symbf{r}=(10,2,2,8)$
\par\smallskip\noindent\textcolor{rulegray}{\rule{\linewidth}{0.3pt}}
\begin{minipage}[c]{0.495\linewidth}\centering\fontsize{9.7}{12}\selectfont
\setlength{\abovedisplayskip}{4pt}\setlength{\belowdisplayskip}{4pt}
\setlength{\abovedisplayshortskip}{4pt}\setlength{\belowdisplayshortskip}{4pt}
\setlength{\arraycolsep}{3pt}\renewcommand{\arraystretch}{1.08}
\[N_{+}^{(12)}(z)=\begin{pmatrix}0 & 0 & 0 & 0\\0 & 0 & 0 & 0\\0 & 0 & 0 & z\\0 & z^{6} + z^{2} & z^{7} + z & 0\end{pmatrix}\]
\end{minipage}%
\begin{minipage}[c]{0.495\linewidth}\centering\fontsize{9.7}{12}\selectfont
\setlength{\abovedisplayskip}{4pt}\setlength{\belowdisplayskip}{4pt}
\setlength{\abovedisplayshortskip}{4pt}\setlength{\belowdisplayshortskip}{4pt}
\setlength{\arraycolsep}{3pt}\renewcommand{\arraystretch}{1.08}
\[N_{-}^{(12)}(z)=\begin{pmatrix}0 & 0 & z^{9} + z & z^{8} + z^{2}\\0 & 0 & z & 0\\z & z & 0 & 0\\0 & z^{4} & 0 & 0\end{pmatrix}\]
\end{minipage}%
\par\noindent
\begin{minipage}[c]{0.495\linewidth}\centering\fontsize{9.7}{12}\selectfont
\setlength{\abovedisplayskip}{4pt}\setlength{\belowdisplayskip}{4pt}
\setlength{\abovedisplayshortskip}{4pt}\setlength{\belowdisplayshortskip}{4pt}
\setlength{\arraycolsep}{3pt}\renewcommand{\arraystretch}{1.08}
\[A_{+}(1)=\begin{pmatrix}2 & 0 & 0 & 0\\0 & 2 & 0 & 0\\0 & 0 & 2 & -1\\0 & -2 & -2 & 2\end{pmatrix}\]
\end{minipage}%
\begin{minipage}[c]{0.495\linewidth}\centering\fontsize{9.7}{12}\selectfont
\setlength{\abovedisplayskip}{4pt}\setlength{\belowdisplayskip}{4pt}
\setlength{\abovedisplayshortskip}{4pt}\setlength{\belowdisplayshortskip}{4pt}
\setlength{\arraycolsep}{3pt}\renewcommand{\arraystretch}{1.08}
\[A_{-}(1)=\begin{pmatrix}2 & 0 & -2 & -2\\0 & 2 & -1 & 0\\-1 & -1 & 2 & 0\\0 & -1 & 0 & 2\end{pmatrix}\]
\end{minipage}%
\par\smallskip\noindent\small $h_+=10,\quad h_-=22,\quad \Omega=32.$\quad Vertices per slice: 11.
\end{minipage}\par\vfill
\noindent\begin{minipage}{\linewidth}
\noindent\textbf{Class 13}\quad\small (constant ID 65)\hfill $\symbf{r}=(2,2,2,4)$
\par\smallskip\noindent\textcolor{rulegray}{\rule{\linewidth}{0.3pt}}
\begin{minipage}[c]{0.495\linewidth}\centering\fontsize{9.7}{12}\selectfont
\setlength{\abovedisplayskip}{4pt}\setlength{\belowdisplayskip}{4pt}
\setlength{\abovedisplayshortskip}{4pt}\setlength{\belowdisplayshortskip}{4pt}
\setlength{\arraycolsep}{3pt}\renewcommand{\arraystretch}{1.08}
\[N_{+}^{(13)}(z)=\begin{pmatrix}0 & 0 & 0 & 0\\0 & 0 & 0 & 0\\0 & 0 & 0 & z\\z^{2} & z^{2} & z^{3} + z & 0\end{pmatrix}\]
\end{minipage}%
\begin{minipage}[c]{0.495\linewidth}\centering\fontsize{9.7}{12}\selectfont
\setlength{\abovedisplayskip}{4pt}\setlength{\belowdisplayskip}{4pt}
\setlength{\abovedisplayshortskip}{4pt}\setlength{\belowdisplayshortskip}{4pt}
\setlength{\arraycolsep}{3pt}\renewcommand{\arraystretch}{1.08}
\[N_{-}^{(13)}(z)=\begin{pmatrix}0 & 0 & z & 0\\0 & 0 & z & 0\\z & z & 0 & 0\\0 & 0 & 0 & 0\end{pmatrix}\]
\end{minipage}%
\par\noindent
\begin{minipage}[c]{0.495\linewidth}\centering\fontsize{9.7}{12}\selectfont
\setlength{\abovedisplayskip}{4pt}\setlength{\belowdisplayskip}{4pt}
\setlength{\abovedisplayshortskip}{4pt}\setlength{\belowdisplayshortskip}{4pt}
\setlength{\arraycolsep}{3pt}\renewcommand{\arraystretch}{1.08}
\[A_{+}(1)=\begin{pmatrix}2 & 0 & 0 & 0\\0 & 2 & 0 & 0\\0 & 0 & 2 & -1\\-1 & -1 & -2 & 2\end{pmatrix}\]
\end{minipage}%
\begin{minipage}[c]{0.495\linewidth}\centering\fontsize{9.7}{12}\selectfont
\setlength{\abovedisplayskip}{4pt}\setlength{\belowdisplayskip}{4pt}
\setlength{\abovedisplayshortskip}{4pt}\setlength{\belowdisplayshortskip}{4pt}
\setlength{\arraycolsep}{3pt}\renewcommand{\arraystretch}{1.08}
\[A_{-}(1)=\begin{pmatrix}2 & 0 & -1 & 0\\0 & 2 & -1 & 0\\-1 & -1 & 2 & 0\\0 & 0 & 0 & 2\end{pmatrix}\]
\end{minipage}%
\par\smallskip\noindent\small $h_+=6,\quad h_-=4,\quad \Omega=20.$\quad Vertices per slice: 5.
\end{minipage}\par\vfill
\noindent\begin{minipage}{\linewidth}
\noindent\textbf{Class 14}\quad\small (constant ID 79)\hfill $\symbf{r}=(2,2,2,10)$
\par\smallskip\noindent\textcolor{rulegray}{\rule{\linewidth}{0.3pt}}
\begin{minipage}[c]{0.495\linewidth}\centering\fontsize{9.7}{12}\selectfont
\setlength{\abovedisplayskip}{4pt}\setlength{\belowdisplayskip}{4pt}
\setlength{\abovedisplayshortskip}{4pt}\setlength{\belowdisplayshortskip}{4pt}
\setlength{\arraycolsep}{3pt}\renewcommand{\arraystretch}{1.08}
\[N_{+}^{(14)}(z)=\begin{pmatrix}0 & 0 & 0 & 0\\0 & 0 & 0 & 0\\0 & 0 & 0 & z\\z^{7} + z^{3} & z^{8} + z^{2} & z^{9} + z & 0\end{pmatrix}\]
\end{minipage}%
\begin{minipage}[c]{0.495\linewidth}\centering\fontsize{9.7}{12}\selectfont
\setlength{\abovedisplayskip}{4pt}\setlength{\belowdisplayskip}{4pt}
\setlength{\abovedisplayshortskip}{4pt}\setlength{\belowdisplayshortskip}{4pt}
\setlength{\arraycolsep}{3pt}\renewcommand{\arraystretch}{1.08}
\[N_{-}^{(14)}(z)=\begin{pmatrix}0 & z & 0 & 0\\z & 0 & z & 0\\0 & z & 0 & 0\\z^{5} & 0 & 0 & 0\end{pmatrix}\]
\end{minipage}%
\par\noindent
\begin{minipage}[c]{0.495\linewidth}\centering\fontsize{9.7}{12}\selectfont
\setlength{\abovedisplayskip}{4pt}\setlength{\belowdisplayskip}{4pt}
\setlength{\abovedisplayshortskip}{4pt}\setlength{\belowdisplayshortskip}{4pt}
\setlength{\arraycolsep}{3pt}\renewcommand{\arraystretch}{1.08}
\[A_{+}(1)=\begin{pmatrix}2 & 0 & 0 & 0\\0 & 2 & 0 & 0\\0 & 0 & 2 & -1\\-2 & -2 & -2 & 2\end{pmatrix}\]
\end{minipage}%
\begin{minipage}[c]{0.495\linewidth}\centering\fontsize{9.7}{12}\selectfont
\setlength{\abovedisplayskip}{4pt}\setlength{\belowdisplayskip}{4pt}
\setlength{\abovedisplayshortskip}{4pt}\setlength{\belowdisplayshortskip}{4pt}
\setlength{\arraycolsep}{3pt}\renewcommand{\arraystretch}{1.08}
\[A_{-}(1)=\begin{pmatrix}2 & -1 & 0 & 0\\-1 & 2 & -1 & 0\\0 & -1 & 2 & 0\\-1 & 0 & 0 & 2\end{pmatrix}\]
\end{minipage}%
\par\smallskip\noindent\small $h_+=12,\quad h_-=10,\quad \Omega=22.$\quad Vertices per slice: 8.
\end{minipage}\par\vfill
\clearpage
\subsection*{Classes 15--17}
\noindent $B_\pm^{(k)}(z)=\operatorname{diag}(1+z^{r_1},\ldots,1+z^{r_4})-N_\pm^{(k)}(z)$.\par\medskip
\noindent\begin{minipage}{\linewidth}
\noindent\textbf{Class 15}\quad\small (constant ID 82)\hfill $\symbf{r}=(2,2,2,7)$
\par\smallskip\noindent\textcolor{rulegray}{\rule{\linewidth}{0.3pt}}
\begin{minipage}[c]{0.495\linewidth}\centering\fontsize{9.7}{12}\selectfont
\setlength{\abovedisplayskip}{4pt}\setlength{\belowdisplayskip}{4pt}
\setlength{\abovedisplayshortskip}{4pt}\setlength{\belowdisplayshortskip}{4pt}
\setlength{\arraycolsep}{3pt}\renewcommand{\arraystretch}{1.08}
\[N_{+}^{(15)}(z)=\begin{pmatrix}0 & 0 & 0 & 0\\0 & 0 & 0 & 0\\0 & 0 & 0 & z\\z^{5} + z^{2} & z^{4} + z^{3} & z^{6} + z & 0\end{pmatrix}\]
\end{minipage}%
\begin{minipage}[c]{0.495\linewidth}\centering\fontsize{9.7}{12}\selectfont
\setlength{\abovedisplayskip}{4pt}\setlength{\belowdisplayskip}{4pt}
\setlength{\abovedisplayshortskip}{4pt}\setlength{\belowdisplayshortskip}{4pt}
\setlength{\arraycolsep}{3pt}\renewcommand{\arraystretch}{1.08}
\[N_{-}^{(15)}(z)=\begin{pmatrix}0 & z & z & 0\\z & z & 0 & 0\\z & 0 & 0 & 0\\0 & 0 & 0 & 0\end{pmatrix}\]
\end{minipage}%
\par\noindent
\begin{minipage}[c]{0.495\linewidth}\centering\fontsize{9.7}{12}\selectfont
\setlength{\abovedisplayskip}{4pt}\setlength{\belowdisplayskip}{4pt}
\setlength{\abovedisplayshortskip}{4pt}\setlength{\belowdisplayshortskip}{4pt}
\setlength{\arraycolsep}{3pt}\renewcommand{\arraystretch}{1.08}
\[A_{+}(1)=\begin{pmatrix}2 & 0 & 0 & 0\\0 & 2 & 0 & 0\\0 & 0 & 2 & -1\\-2 & -2 & -2 & 2\end{pmatrix}\]
\end{minipage}%
\begin{minipage}[c]{0.495\linewidth}\centering\fontsize{9.7}{12}\selectfont
\setlength{\abovedisplayskip}{4pt}\setlength{\belowdisplayskip}{4pt}
\setlength{\abovedisplayshortskip}{4pt}\setlength{\belowdisplayshortskip}{4pt}
\setlength{\arraycolsep}{3pt}\renewcommand{\arraystretch}{1.08}
\[A_{-}(1)=\begin{pmatrix}2 & -1 & -1 & 0\\-1 & 1 & 0 & 0\\-1 & 0 & 2 & 0\\0 & 0 & 0 & 2\end{pmatrix}\]
\end{minipage}%
\par\smallskip\noindent\small $h_+=9,\quad h_-=7,\quad \Omega=16.$\quad Vertices per slice: 13.
\end{minipage}\par\vfill
\noindent\begin{minipage}{\linewidth}
\noindent\textbf{Class 16}\quad\small (constant ID 122)\hfill $\symbf{r}=(3,3,2,2)$
\par\smallskip\noindent\textcolor{rulegray}{\rule{\linewidth}{0.3pt}}
\begin{minipage}[c]{0.495\linewidth}\centering\fontsize{9.7}{12}\selectfont
\setlength{\abovedisplayskip}{4pt}\setlength{\belowdisplayskip}{4pt}
\setlength{\abovedisplayshortskip}{4pt}\setlength{\belowdisplayshortskip}{4pt}
\setlength{\arraycolsep}{3pt}\renewcommand{\arraystretch}{1.08}
\[N_{+}^{(16)}(z)=\begin{pmatrix}0 & 0 & 0 & 0\\0 & 0 & 0 & 0\\0 & 0 & z & 0\\0 & 0 & 0 & z\end{pmatrix}\]
\end{minipage}%
\begin{minipage}[c]{0.495\linewidth}\centering\fontsize{9.7}{12}\selectfont
\setlength{\abovedisplayskip}{4pt}\setlength{\belowdisplayskip}{4pt}
\setlength{\abovedisplayshortskip}{4pt}\setlength{\belowdisplayshortskip}{4pt}
\setlength{\arraycolsep}{3pt}\renewcommand{\arraystretch}{1.08}
\[N_{-}^{(16)}(z)=\begin{pmatrix}0 & z & 0 & 0\\z^{2} & 0 & 0 & z^{2} + z\\0 & 0 & 0 & z\\0 & z & z & 0\end{pmatrix}\]
\end{minipage}%
\par\noindent
\begin{minipage}[c]{0.495\linewidth}\centering\fontsize{9.7}{12}\selectfont
\setlength{\abovedisplayskip}{4pt}\setlength{\belowdisplayskip}{4pt}
\setlength{\abovedisplayshortskip}{4pt}\setlength{\belowdisplayshortskip}{4pt}
\setlength{\arraycolsep}{3pt}\renewcommand{\arraystretch}{1.08}
\[A_{+}(1)=\begin{pmatrix}2 & 0 & 0 & 0\\0 & 2 & 0 & 0\\0 & 0 & 1 & 0\\0 & 0 & 0 & 1\end{pmatrix}\]
\end{minipage}%
\begin{minipage}[c]{0.495\linewidth}\centering\fontsize{9.7}{12}\selectfont
\setlength{\abovedisplayskip}{4pt}\setlength{\belowdisplayskip}{4pt}
\setlength{\abovedisplayshortskip}{4pt}\setlength{\belowdisplayshortskip}{4pt}
\setlength{\arraycolsep}{3pt}\renewcommand{\arraystretch}{1.08}
\[A_{-}(1)=\begin{pmatrix}2 & -1 & 0 & 0\\-1 & 2 & 0 & -2\\0 & 0 & 2 & -1\\0 & -1 & -1 & 2\end{pmatrix}\]
\end{minipage}%
\par\smallskip\noindent\small $h_+=3,\quad h_-=15,\quad \Omega=18.$\quad Vertices per slice: 10.
\end{minipage}\par\vfill
\noindent\begin{minipage}{\linewidth}
\noindent\textbf{Class 17}\quad\small (constant ID 136)\hfill $\symbf{r}=(2,2,8,8)$
\par\smallskip\noindent\textcolor{rulegray}{\rule{\linewidth}{0.3pt}}
\begin{minipage}[c]{0.495\linewidth}\centering\fontsize{9.7}{12}\selectfont
\setlength{\abovedisplayskip}{4pt}\setlength{\belowdisplayskip}{4pt}
\setlength{\abovedisplayshortskip}{4pt}\setlength{\belowdisplayshortskip}{4pt}
\setlength{\arraycolsep}{3pt}\renewcommand{\arraystretch}{1.08}
\[N_{+}^{(17)}(z)=\begin{pmatrix}0 & 0 & 0 & 0\\0 & 0 & 0 & 0\\0 & 0 & z^{4} & 0\\z^{4} & z^{5} + z^{3} & 0 & z^{4}\end{pmatrix}\]
\end{minipage}%
\begin{minipage}[c]{0.495\linewidth}\centering\fontsize{9.7}{12}\selectfont
\setlength{\abovedisplayskip}{4pt}\setlength{\belowdisplayskip}{4pt}
\setlength{\abovedisplayshortskip}{4pt}\setlength{\belowdisplayshortskip}{4pt}
\setlength{\arraycolsep}{3pt}\renewcommand{\arraystretch}{1.08}
\[N_{-}^{(17)}(z)=\begin{pmatrix}0 & z & 0 & 0\\z & 0 & 0 & z\\0 & 0 & 0 & z\\0 & z^{7} + z & z^{7} & 0\end{pmatrix}\]
\end{minipage}%
\par\noindent
\begin{minipage}[c]{0.495\linewidth}\centering\fontsize{9.7}{12}\selectfont
\setlength{\abovedisplayskip}{4pt}\setlength{\belowdisplayskip}{4pt}
\setlength{\abovedisplayshortskip}{4pt}\setlength{\belowdisplayshortskip}{4pt}
\setlength{\arraycolsep}{3pt}\renewcommand{\arraystretch}{1.08}
\[A_{+}(1)=\begin{pmatrix}2 & 0 & 0 & 0\\0 & 2 & 0 & 0\\0 & 0 & 1 & 0\\-1 & -2 & 0 & 1\end{pmatrix}\]
\end{minipage}%
\begin{minipage}[c]{0.495\linewidth}\centering\fontsize{9.7}{12}\selectfont
\setlength{\abovedisplayskip}{4pt}\setlength{\belowdisplayskip}{4pt}
\setlength{\abovedisplayshortskip}{4pt}\setlength{\belowdisplayshortskip}{4pt}
\setlength{\arraycolsep}{3pt}\renewcommand{\arraystretch}{1.08}
\[A_{-}(1)=\begin{pmatrix}2 & -1 & 0 & 0\\-1 & 2 & 0 & -1\\0 & 0 & 2 & -1\\0 & -2 & -1 & 2\end{pmatrix}\]
\end{minipage}%
\par\smallskip\noindent\small $h_+=12,\quad h_-=30,\quad \Omega=42.$\quad Vertices per slice: 10.
\end{minipage}\par\vfill
\clearpage
\subsection*{Classes 18--20}
\noindent $B_\pm^{(k)}(z)=\operatorname{diag}(1+z^{r_1},\ldots,1+z^{r_4})-N_\pm^{(k)}(z)$.\par\medskip
\noindent\begin{minipage}{\linewidth}
\noindent\textbf{Class 18}\quad\small (constant ID 207)\hfill $\symbf{r}=(14,2,6,8)$
\par\smallskip\noindent\textcolor{rulegray}{\rule{\linewidth}{0.3pt}}
\begin{minipage}[c]{0.495\linewidth}\centering\fontsize{9.7}{12}\selectfont
\setlength{\abovedisplayskip}{4pt}\setlength{\belowdisplayskip}{4pt}
\setlength{\abovedisplayshortskip}{4pt}\setlength{\belowdisplayshortskip}{4pt}
\setlength{\arraycolsep}{3pt}\renewcommand{\arraystretch}{1.08}
\[N_{+}^{(18)}(z)=\begin{pmatrix}0 & 0 & 0 & 0\\0 & 0 & 0 & 0\\0 & z^{3} & 0 & z^{5}\\0 & z^{2} & z^{3} + z & 0\end{pmatrix}\]
\end{minipage}%
\begin{minipage}[c]{0.495\linewidth}\centering\fontsize{9.7}{12}\selectfont
\setlength{\abovedisplayskip}{4pt}\setlength{\belowdisplayskip}{4pt}
\setlength{\abovedisplayshortskip}{4pt}\setlength{\belowdisplayshortskip}{4pt}
\setlength{\arraycolsep}{3pt}\renewcommand{\arraystretch}{1.08}
\[N_{-}^{(18)}(z)=\begin{pmatrix}0 & 0 & z^{2} & z^{7} + z\\0 & 0 & z & 0\\0 & z^{5} + z & 0 & 0\\z^{7} & 0 & 0 & 0\end{pmatrix}\]
\end{minipage}%
\par\noindent
\begin{minipage}[c]{0.495\linewidth}\centering\fontsize{9.7}{12}\selectfont
\setlength{\abovedisplayskip}{4pt}\setlength{\belowdisplayskip}{4pt}
\setlength{\abovedisplayshortskip}{4pt}\setlength{\belowdisplayshortskip}{4pt}
\setlength{\arraycolsep}{3pt}\renewcommand{\arraystretch}{1.08}
\[A_{+}(1)=\begin{pmatrix}2 & 0 & 0 & 0\\0 & 2 & 0 & 0\\0 & -1 & 2 & -1\\0 & -1 & -2 & 2\end{pmatrix}\]
\end{minipage}%
\begin{minipage}[c]{0.495\linewidth}\centering\fontsize{9.7}{12}\selectfont
\setlength{\abovedisplayskip}{4pt}\setlength{\belowdisplayskip}{4pt}
\setlength{\abovedisplayshortskip}{4pt}\setlength{\belowdisplayshortskip}{4pt}
\setlength{\arraycolsep}{3pt}\renewcommand{\arraystretch}{1.08}
\[A_{-}(1)=\begin{pmatrix}2 & 0 & -1 & -2\\0 & 2 & -1 & 0\\0 & -2 & 2 & 0\\-1 & 0 & 0 & 2\end{pmatrix}\]
\end{minipage}%
\par\smallskip\noindent\small $h_+=14,\quad h_-=22,\quad \Omega=36.$\quad Vertices per slice: 15.
\end{minipage}\par\vfill
\noindent\begin{minipage}{\linewidth}
\noindent\textbf{Class 19}\quad\small (constant ID 413)\hfill $\symbf{r}=(2,2,16,12)$
\par\smallskip\noindent\textcolor{rulegray}{\rule{\linewidth}{0.3pt}}
\begin{minipage}[c]{0.495\linewidth}\centering\fontsize{9.7}{12}\selectfont
\setlength{\abovedisplayskip}{4pt}\setlength{\belowdisplayskip}{4pt}
\setlength{\abovedisplayshortskip}{4pt}\setlength{\belowdisplayshortskip}{4pt}
\setlength{\arraycolsep}{3pt}\renewcommand{\arraystretch}{1.08}
\[N_{+}^{(19)}(z)=\begin{pmatrix}0 & 0 & 0 & 0\\0 & 0 & 0 & 0\\z^{3} & z^{4} + z^{2} & 0 & z^{5} + z^{3} + z\\z^{8} + z^{4} & z^{9} + z^{3} & z^{11} & 0\end{pmatrix}\]
\end{minipage}%
\begin{minipage}[c]{0.495\linewidth}\centering\fontsize{9.7}{12}\selectfont
\setlength{\abovedisplayskip}{4pt}\setlength{\belowdisplayskip}{4pt}
\setlength{\abovedisplayshortskip}{4pt}\setlength{\belowdisplayshortskip}{4pt}
\setlength{\arraycolsep}{3pt}\renewcommand{\arraystretch}{1.08}
\[N_{-}^{(19)}(z)=\begin{pmatrix}0 & z & 0 & 0\\z & 0 & 0 & z\\0 & 0 & z^{8} & 0\\z^{6} & z^{11} + z & 0 & 0\end{pmatrix}\]
\end{minipage}%
\par\noindent
\begin{minipage}[c]{0.495\linewidth}\centering\fontsize{9.7}{12}\selectfont
\setlength{\abovedisplayskip}{4pt}\setlength{\belowdisplayskip}{4pt}
\setlength{\abovedisplayshortskip}{4pt}\setlength{\belowdisplayshortskip}{4pt}
\setlength{\arraycolsep}{3pt}\renewcommand{\arraystretch}{1.08}
\[A_{+}(1)=\begin{pmatrix}2 & 0 & 0 & 0\\0 & 2 & 0 & 0\\-1 & -2 & 2 & -3\\-2 & -2 & -1 & 2\end{pmatrix}\]
\end{minipage}%
\begin{minipage}[c]{0.495\linewidth}\centering\fontsize{9.7}{12}\selectfont
\setlength{\abovedisplayskip}{4pt}\setlength{\belowdisplayskip}{4pt}
\setlength{\abovedisplayshortskip}{4pt}\setlength{\belowdisplayshortskip}{4pt}
\setlength{\arraycolsep}{3pt}\renewcommand{\arraystretch}{1.08}
\[A_{-}(1)=\begin{pmatrix}2 & -1 & 0 & 0\\-1 & 2 & 0 & -1\\0 & 0 & 1 & 0\\-1 & -2 & 0 & 2\end{pmatrix}\]
\end{minipage}%
\par\smallskip\noindent\small $h_+=42,\quad h_-=24,\quad \Omega=66.$\quad Vertices per slice: 16.
\end{minipage}\par\vfill
\noindent\begin{minipage}{\linewidth}
\noindent\textbf{Class 20}\quad\small (constant ID 432)\hfill $\symbf{r}=(8,8,2,2)$
\par\smallskip\noindent\textcolor{rulegray}{\rule{\linewidth}{0.3pt}}
\begin{minipage}[c]{0.495\linewidth}\centering\fontsize{9.7}{12}\selectfont
\setlength{\abovedisplayskip}{4pt}\setlength{\belowdisplayskip}{4pt}
\setlength{\abovedisplayshortskip}{4pt}\setlength{\belowdisplayshortskip}{4pt}
\setlength{\arraycolsep}{3pt}\renewcommand{\arraystretch}{1.08}
\[N_{+}^{(20)}(z)=\begin{pmatrix}0 & 0 & 0 & 0\\0 & 0 & 0 & z^{4}\\0 & 0 & 0 & z\\0 & 0 & z & 0\end{pmatrix}\]
\end{minipage}%
\begin{minipage}[c]{0.495\linewidth}\centering\fontsize{9.7}{12}\selectfont
\setlength{\abovedisplayskip}{4pt}\setlength{\belowdisplayskip}{4pt}
\setlength{\abovedisplayshortskip}{4pt}\setlength{\belowdisplayshortskip}{4pt}
\setlength{\arraycolsep}{3pt}\renewcommand{\arraystretch}{1.08}
\[N_{-}^{(20)}(z)=\begin{pmatrix}0 & z & 0 & 0\\z^{7} & 0 & z^{7} + z & z^{6} + z^{2}\\0 & z & 0 & 0\\0 & 0 & 0 & 0\end{pmatrix}\]
\end{minipage}%
\par\noindent
\begin{minipage}[c]{0.495\linewidth}\centering\fontsize{9.7}{12}\selectfont
\setlength{\abovedisplayskip}{4pt}\setlength{\belowdisplayskip}{4pt}
\setlength{\abovedisplayshortskip}{4pt}\setlength{\belowdisplayshortskip}{4pt}
\setlength{\arraycolsep}{3pt}\renewcommand{\arraystretch}{1.08}
\[A_{+}(1)=\begin{pmatrix}2 & 0 & 0 & 0\\0 & 2 & 0 & -1\\0 & 0 & 2 & -1\\0 & 0 & -1 & 2\end{pmatrix}\]
\end{minipage}%
\begin{minipage}[c]{0.495\linewidth}\centering\fontsize{9.7}{12}\selectfont
\setlength{\abovedisplayskip}{4pt}\setlength{\belowdisplayskip}{4pt}
\setlength{\abovedisplayshortskip}{4pt}\setlength{\belowdisplayshortskip}{4pt}
\setlength{\arraycolsep}{3pt}\renewcommand{\arraystretch}{1.08}
\[A_{-}(1)=\begin{pmatrix}2 & -1 & 0 & 0\\-1 & 2 & -2 & -2\\0 & -1 & 2 & 0\\0 & 0 & 0 & 2\end{pmatrix}\]
\end{minipage}%
\par\smallskip\noindent\small $h_+=8,\quad h_-=18,\quad \Omega=26.$\quad Vertices per slice: 10.
\end{minipage}\par\vfill
\clearpage
\subsection*{Classes 21--23}
\noindent $B_\pm^{(k)}(z)=\operatorname{diag}(1+z^{r_1},\ldots,1+z^{r_4})-N_\pm^{(k)}(z)$.\par\medskip
\noindent\begin{minipage}{\linewidth}
\noindent\textbf{Class 21}\quad\small (constant ID 436)\hfill $\symbf{r}=(14,6,2,6)$
\par\smallskip\noindent\textcolor{rulegray}{\rule{\linewidth}{0.3pt}}
\begin{minipage}[c]{0.495\linewidth}\centering\fontsize{9.7}{12}\selectfont
\setlength{\abovedisplayskip}{4pt}\setlength{\belowdisplayskip}{4pt}
\setlength{\abovedisplayshortskip}{4pt}\setlength{\belowdisplayshortskip}{4pt}
\setlength{\arraycolsep}{3pt}\renewcommand{\arraystretch}{1.08}
\[N_{+}^{(21)}(z)=\begin{pmatrix}0 & 0 & 0 & 0\\0 & 0 & 0 & z\\0 & 0 & 0 & z\\0 & z^{5} & z^{5} + z & 0\end{pmatrix}\]
\end{minipage}%
\begin{minipage}[c]{0.495\linewidth}\centering\fontsize{9.7}{12}\selectfont
\setlength{\abovedisplayskip}{4pt}\setlength{\belowdisplayskip}{4pt}
\setlength{\abovedisplayshortskip}{4pt}\setlength{\belowdisplayshortskip}{4pt}
\setlength{\arraycolsep}{3pt}\renewcommand{\arraystretch}{1.08}
\[N_{-}^{(21)}(z)=\begin{pmatrix}0 & z^{9} + z & z^{3} & z^{4} + z^{2}\\z^{5} & 0 & 0 & 0\\0 & 0 & 0 & 0\\0 & 0 & z^{3} & 0\end{pmatrix}\]
\end{minipage}%
\par\noindent
\begin{minipage}[c]{0.495\linewidth}\centering\fontsize{9.7}{12}\selectfont
\setlength{\abovedisplayskip}{4pt}\setlength{\belowdisplayskip}{4pt}
\setlength{\abovedisplayshortskip}{4pt}\setlength{\belowdisplayshortskip}{4pt}
\setlength{\arraycolsep}{3pt}\renewcommand{\arraystretch}{1.08}
\[A_{+}(1)=\begin{pmatrix}2 & 0 & 0 & 0\\0 & 2 & 0 & -1\\0 & 0 & 2 & -1\\0 & -1 & -2 & 2\end{pmatrix}\]
\end{minipage}%
\begin{minipage}[c]{0.495\linewidth}\centering\fontsize{9.7}{12}\selectfont
\setlength{\abovedisplayskip}{4pt}\setlength{\belowdisplayskip}{4pt}
\setlength{\abovedisplayshortskip}{4pt}\setlength{\belowdisplayshortskip}{4pt}
\setlength{\arraycolsep}{3pt}\renewcommand{\arraystretch}{1.08}
\[A_{-}(1)=\begin{pmatrix}2 & -2 & -1 & -2\\-1 & 2 & 0 & 0\\0 & 0 & 2 & 0\\0 & 0 & -1 & 2\end{pmatrix}\]
\end{minipage}%
\par\smallskip\noindent\small $h_+=14,\quad h_-=20,\quad \Omega=34.$\quad Vertices per slice: 14.
\end{minipage}\par\vfill
\noindent\begin{minipage}{\linewidth}
\noindent\textbf{Class 22}\quad\small (constant ID 437)\hfill $\symbf{r}=(8,3,2,3)$
\par\smallskip\noindent\textcolor{rulegray}{\rule{\linewidth}{0.3pt}}
\begin{minipage}[c]{0.495\linewidth}\centering\fontsize{9.7}{12}\selectfont
\setlength{\abovedisplayskip}{4pt}\setlength{\belowdisplayskip}{4pt}
\setlength{\abovedisplayshortskip}{4pt}\setlength{\belowdisplayshortskip}{4pt}
\setlength{\arraycolsep}{3pt}\renewcommand{\arraystretch}{1.08}
\[N_{+}^{(22)}(z)=\begin{pmatrix}0 & 0 & 0 & 0\\0 & 0 & 0 & z\\0 & 0 & 0 & z\\0 & z^{2} & z^{2} + z & 0\end{pmatrix}\]
\end{minipage}%
\begin{minipage}[c]{0.495\linewidth}\centering\fontsize{9.7}{12}\selectfont
\setlength{\abovedisplayskip}{4pt}\setlength{\belowdisplayskip}{4pt}
\setlength{\abovedisplayshortskip}{4pt}\setlength{\belowdisplayshortskip}{4pt}
\setlength{\arraycolsep}{3pt}\renewcommand{\arraystretch}{1.08}
\[N_{-}^{(22)}(z)=\begin{pmatrix}0 & z^{6} + z & z^{3} & z^{4} + z^{2}\\z^{2} & 0 & 0 & 0\\0 & 0 & z & 0\\0 & 0 & 0 & 0\end{pmatrix}\]
\end{minipage}%
\par\noindent
\begin{minipage}[c]{0.495\linewidth}\centering\fontsize{9.7}{12}\selectfont
\setlength{\abovedisplayskip}{4pt}\setlength{\belowdisplayskip}{4pt}
\setlength{\abovedisplayshortskip}{4pt}\setlength{\belowdisplayshortskip}{4pt}
\setlength{\arraycolsep}{3pt}\renewcommand{\arraystretch}{1.08}
\[A_{+}(1)=\begin{pmatrix}2 & 0 & 0 & 0\\0 & 2 & 0 & -1\\0 & 0 & 2 & -1\\0 & -1 & -2 & 2\end{pmatrix}\]
\end{minipage}%
\begin{minipage}[c]{0.495\linewidth}\centering\fontsize{9.7}{12}\selectfont
\setlength{\abovedisplayskip}{4pt}\setlength{\belowdisplayskip}{4pt}
\setlength{\abovedisplayshortskip}{4pt}\setlength{\belowdisplayshortskip}{4pt}
\setlength{\arraycolsep}{3pt}\renewcommand{\arraystretch}{1.08}
\[A_{-}(1)=\begin{pmatrix}2 & -2 & -1 & -2\\-1 & 2 & 0 & 0\\0 & 0 & 1 & 0\\0 & 0 & 0 & 2\end{pmatrix}\]
\end{minipage}%
\par\smallskip\noindent\small $h_+=8,\quad h_-=11,\quad \Omega=19.$\quad Vertices per slice: 16.
\end{minipage}\par\vfill
\noindent\begin{minipage}{\linewidth}
\noindent\textbf{Class 23}\quad\small (constant ID 453)\hfill $\symbf{r}=(2,12,2,12)$
\par\smallskip\noindent\textcolor{rulegray}{\rule{\linewidth}{0.3pt}}
\begin{minipage}[c]{0.495\linewidth}\centering\fontsize{9.7}{12}\selectfont
\setlength{\abovedisplayskip}{4pt}\setlength{\belowdisplayskip}{4pt}
\setlength{\abovedisplayshortskip}{4pt}\setlength{\belowdisplayshortskip}{4pt}
\setlength{\arraycolsep}{3pt}\renewcommand{\arraystretch}{1.08}
\[N_{+}^{(23)}(z)=\begin{pmatrix}0 & 0 & 0 & 0\\0 & 0 & 0 & z^{11}\\0 & 0 & 0 & z\\z^{10} + z^{6} + z^{2} & z & z^{11} + z & 0\end{pmatrix}\]
\end{minipage}%
\begin{minipage}[c]{0.495\linewidth}\centering\fontsize{9.7}{12}\selectfont
\setlength{\abovedisplayskip}{4pt}\setlength{\belowdisplayskip}{4pt}
\setlength{\abovedisplayshortskip}{4pt}\setlength{\belowdisplayshortskip}{4pt}
\setlength{\arraycolsep}{3pt}\renewcommand{\arraystretch}{1.08}
\[N_{-}^{(23)}(z)=\begin{pmatrix}0 & 0 & z & 0\\0 & z^{6} & 0 & 0\\z & 0 & 0 & 0\\z^{8} + z^{4} & 0 & 0 & z^{6}\end{pmatrix}\]
\end{minipage}%
\par\noindent
\begin{minipage}[c]{0.495\linewidth}\centering\fontsize{9.7}{12}\selectfont
\setlength{\abovedisplayskip}{4pt}\setlength{\belowdisplayskip}{4pt}
\setlength{\abovedisplayshortskip}{4pt}\setlength{\belowdisplayshortskip}{4pt}
\setlength{\arraycolsep}{3pt}\renewcommand{\arraystretch}{1.08}
\[A_{+}(1)=\begin{pmatrix}2 & 0 & 0 & 0\\0 & 2 & 0 & -1\\0 & 0 & 2 & -1\\-3 & -1 & -2 & 2\end{pmatrix}\]
\end{minipage}%
\begin{minipage}[c]{0.495\linewidth}\centering\fontsize{9.7}{12}\selectfont
\setlength{\abovedisplayskip}{4pt}\setlength{\belowdisplayskip}{4pt}
\setlength{\abovedisplayshortskip}{4pt}\setlength{\belowdisplayshortskip}{4pt}
\setlength{\arraycolsep}{3pt}\renewcommand{\arraystretch}{1.08}
\[A_{-}(1)=\begin{pmatrix}2 & 0 & -1 & 0\\0 & 1 & 0 & 0\\-1 & 0 & 2 & 0\\-2 & 0 & 0 & 1\end{pmatrix}\]
\end{minipage}%
\par\smallskip\noindent\small $h_+=26,\quad h_-=18,\quad \Omega=44.$\quad Vertices per slice: 14.
\end{minipage}\par\vfill
\clearpage
\subsection*{Classes 24--26}
\noindent $B_\pm^{(k)}(z)=\operatorname{diag}(1+z^{r_1},\ldots,1+z^{r_4})-N_\pm^{(k)}(z)$.\par\medskip
\noindent\begin{minipage}{\linewidth}
\noindent\textbf{Class 24}\quad\small (constant ID 456)\hfill $\symbf{r}=(7,2,3,2)$
\par\smallskip\noindent\textcolor{rulegray}{\rule{\linewidth}{0.3pt}}
\begin{minipage}[c]{0.495\linewidth}\centering\fontsize{9.7}{12}\selectfont
\setlength{\abovedisplayskip}{4pt}\setlength{\belowdisplayskip}{4pt}
\setlength{\abovedisplayshortskip}{4pt}\setlength{\belowdisplayshortskip}{4pt}
\setlength{\arraycolsep}{3pt}\renewcommand{\arraystretch}{1.08}
\[N_{+}^{(24)}(z)=\begin{pmatrix}0 & 0 & 0 & 0\\0 & 0 & 0 & z\\0 & 0 & 0 & z^{2} + z\\0 & z & z & 0\end{pmatrix}\]
\end{minipage}%
\begin{minipage}[c]{0.495\linewidth}\centering\fontsize{9.7}{12}\selectfont
\setlength{\abovedisplayskip}{4pt}\setlength{\belowdisplayskip}{4pt}
\setlength{\abovedisplayshortskip}{4pt}\setlength{\belowdisplayshortskip}{4pt}
\setlength{\arraycolsep}{3pt}\renewcommand{\arraystretch}{1.08}
\[N_{-}^{(24)}(z)=\begin{pmatrix}0 & z^{3} & z^{5} + z^{3} + z & z^{4} + z^{2}\\0 & z & 0 & 0\\z^{2} & 0 & 0 & 0\\0 & 0 & 0 & z\end{pmatrix}\]
\end{minipage}%
\par\noindent
\begin{minipage}[c]{0.495\linewidth}\centering\fontsize{9.7}{12}\selectfont
\setlength{\abovedisplayskip}{4pt}\setlength{\belowdisplayskip}{4pt}
\setlength{\abovedisplayshortskip}{4pt}\setlength{\belowdisplayshortskip}{4pt}
\setlength{\arraycolsep}{3pt}\renewcommand{\arraystretch}{1.08}
\[A_{+}(1)=\begin{pmatrix}2 & 0 & 0 & 0\\0 & 2 & 0 & -1\\0 & 0 & 2 & -2\\0 & -1 & -1 & 2\end{pmatrix}\]
\end{minipage}%
\begin{minipage}[c]{0.495\linewidth}\centering\fontsize{9.7}{12}\selectfont
\setlength{\abovedisplayskip}{4pt}\setlength{\belowdisplayskip}{4pt}
\setlength{\abovedisplayshortskip}{4pt}\setlength{\belowdisplayshortskip}{4pt}
\setlength{\arraycolsep}{3pt}\renewcommand{\arraystretch}{1.08}
\[A_{-}(1)=\begin{pmatrix}2 & -1 & -3 & -2\\0 & 1 & 0 & 0\\-1 & 0 & 2 & 0\\0 & 0 & 0 & 1\end{pmatrix}\]
\end{minipage}%
\par\smallskip\noindent\small $h_+=7,\quad h_-=15,\quad \Omega=22.$\quad Vertices per slice: 14.
\end{minipage}\par\vfill
\noindent\begin{minipage}{\linewidth}
\noindent\textbf{Class 25}\quad\small (constant ID 469)\hfill $\symbf{r}=(8,3,2,5)$
\par\smallskip\noindent\textcolor{rulegray}{\rule{\linewidth}{0.3pt}}
\begin{minipage}[c]{0.495\linewidth}\centering\fontsize{9.7}{12}\selectfont
\setlength{\abovedisplayskip}{4pt}\setlength{\belowdisplayskip}{4pt}
\setlength{\abovedisplayshortskip}{4pt}\setlength{\belowdisplayshortskip}{4pt}
\setlength{\arraycolsep}{3pt}\renewcommand{\arraystretch}{1.08}
\[N_{+}^{(25)}(z)=\begin{pmatrix}0 & 0 & 0 & 0\\0 & 0 & 0 & z^{2}\\0 & 0 & z & 0\\0 & z^{3} + z & z^{2} & 0\end{pmatrix}\]
\end{minipage}%
\begin{minipage}[c]{0.495\linewidth}\centering\fontsize{9.7}{12}\selectfont
\setlength{\abovedisplayskip}{4pt}\setlength{\belowdisplayskip}{4pt}
\setlength{\abovedisplayshortskip}{4pt}\setlength{\belowdisplayshortskip}{4pt}
\setlength{\arraycolsep}{3pt}\renewcommand{\arraystretch}{1.08}
\[N_{-}^{(25)}(z)=\begin{pmatrix}0 & z^{2} & 0 & z^{4} + z\\0 & 0 & z^{2} + z & 0\\0 & z & 0 & 0\\z^{4} & 0 & 0 & 0\end{pmatrix}\]
\end{minipage}%
\par\noindent
\begin{minipage}[c]{0.495\linewidth}\centering\fontsize{9.7}{12}\selectfont
\setlength{\abovedisplayskip}{4pt}\setlength{\belowdisplayskip}{4pt}
\setlength{\abovedisplayshortskip}{4pt}\setlength{\belowdisplayshortskip}{4pt}
\setlength{\arraycolsep}{3pt}\renewcommand{\arraystretch}{1.08}
\[A_{+}(1)=\begin{pmatrix}2 & 0 & 0 & 0\\0 & 2 & 0 & -1\\0 & 0 & 1 & 0\\0 & -2 & -1 & 2\end{pmatrix}\]
\end{minipage}%
\begin{minipage}[c]{0.495\linewidth}\centering\fontsize{9.7}{12}\selectfont
\setlength{\abovedisplayskip}{4pt}\setlength{\belowdisplayskip}{4pt}
\setlength{\abovedisplayshortskip}{4pt}\setlength{\belowdisplayshortskip}{4pt}
\setlength{\arraycolsep}{3pt}\renewcommand{\arraystretch}{1.08}
\[A_{-}(1)=\begin{pmatrix}2 & -1 & 0 & -2\\0 & 2 & -2 & 0\\0 & -1 & 2 & 0\\-1 & 0 & 0 & 2\end{pmatrix}\]
\end{minipage}%
\par\smallskip\noindent\small $h_+=8,\quad h_-=13,\quad \Omega=21.$\quad Vertices per slice: 18.
\end{minipage}\par\vfill
\noindent\begin{minipage}{\linewidth}
\noindent\textbf{Class 26}\quad\small (constant ID 949)\hfill $\symbf{r}=(2,12,10,12)$
\par\smallskip\noindent\textcolor{rulegray}{\rule{\linewidth}{0.3pt}}
\begin{minipage}[c]{0.495\linewidth}\centering\fontsize{9.7}{12}\selectfont
\setlength{\abovedisplayskip}{4pt}\setlength{\belowdisplayskip}{4pt}
\setlength{\abovedisplayshortskip}{4pt}\setlength{\belowdisplayshortskip}{4pt}
\setlength{\arraycolsep}{3pt}\renewcommand{\arraystretch}{1.08}
\[N_{+}^{(26)}(z)=\begin{pmatrix}0 & 0 & 0 & 0\\0 & 0 & 0 & z^{11}\\z^{7} + z^{3} & 0 & 0 & z^{9}\\z^{2} & z & z^{3} + z & 0\end{pmatrix}\]
\end{minipage}%
\begin{minipage}[c]{0.495\linewidth}\centering\fontsize{9.7}{12}\selectfont
\setlength{\abovedisplayskip}{4pt}\setlength{\belowdisplayskip}{4pt}
\setlength{\abovedisplayshortskip}{4pt}\setlength{\belowdisplayshortskip}{4pt}
\setlength{\arraycolsep}{3pt}\renewcommand{\arraystretch}{1.08}
\[N_{-}^{(26)}(z)=\begin{pmatrix}0 & 0 & z & 0\\0 & z^{6} & 0 & 0\\z^{9} + z^{5} + z & 0 & 0 & 0\\0 & 0 & 0 & z^{6}\end{pmatrix}\]
\end{minipage}%
\par\noindent
\begin{minipage}[c]{0.495\linewidth}\centering\fontsize{9.7}{12}\selectfont
\setlength{\abovedisplayskip}{4pt}\setlength{\belowdisplayskip}{4pt}
\setlength{\abovedisplayshortskip}{4pt}\setlength{\belowdisplayshortskip}{4pt}
\setlength{\arraycolsep}{3pt}\renewcommand{\arraystretch}{1.08}
\[A_{+}(1)=\begin{pmatrix}2 & 0 & 0 & 0\\0 & 2 & 0 & -1\\-2 & 0 & 2 & -1\\-1 & -1 & -2 & 2\end{pmatrix}\]
\end{minipage}%
\begin{minipage}[c]{0.495\linewidth}\centering\fontsize{9.7}{12}\selectfont
\setlength{\abovedisplayskip}{4pt}\setlength{\belowdisplayskip}{4pt}
\setlength{\abovedisplayshortskip}{4pt}\setlength{\belowdisplayshortskip}{4pt}
\setlength{\arraycolsep}{3pt}\renewcommand{\arraystretch}{1.08}
\[A_{-}(1)=\begin{pmatrix}2 & 0 & -1 & 0\\0 & 1 & 0 & 0\\-3 & 0 & 2 & 0\\0 & 0 & 0 & 1\end{pmatrix}\]
\end{minipage}%
\par\smallskip\noindent\small $h_+=34,\quad h_-=18,\quad \Omega=52.$\quad Vertices per slice: 18.
\end{minipage}\par\vfill
\clearpage
\subsection*{Classes 27--29}
\noindent $B_\pm^{(k)}(z)=\operatorname{diag}(1+z^{r_1},\ldots,1+z^{r_4})-N_\pm^{(k)}(z)$.\par\medskip
\noindent\begin{minipage}{\linewidth}
\noindent\textbf{Class 27}\quad\small (constant ID 1262)\hfill $\symbf{r}=(3,2,2,2)$
\par\smallskip\noindent\textcolor{rulegray}{\rule{\linewidth}{0.3pt}}
\begin{minipage}[c]{0.495\linewidth}\centering\fontsize{9.7}{12}\selectfont
\setlength{\abovedisplayskip}{4pt}\setlength{\belowdisplayskip}{4pt}
\setlength{\abovedisplayshortskip}{4pt}\setlength{\belowdisplayshortskip}{4pt}
\setlength{\arraycolsep}{3pt}\renewcommand{\arraystretch}{1.08}
\[N_{+}^{(27)}(z)=\begin{pmatrix}0 & 0 & 0 & 0\\0 & z & 0 & 0\\0 & 0 & z & 0\\0 & 0 & 0 & z\end{pmatrix}\]
\end{minipage}%
\begin{minipage}[c]{0.495\linewidth}\centering\fontsize{9.7}{12}\selectfont
\setlength{\abovedisplayskip}{4pt}\setlength{\belowdisplayskip}{4pt}
\setlength{\abovedisplayshortskip}{4pt}\setlength{\belowdisplayshortskip}{4pt}
\setlength{\arraycolsep}{3pt}\renewcommand{\arraystretch}{1.08}
\[N_{-}^{(27)}(z)=\begin{pmatrix}0 & 0 & 0 & z^{2} + z\\0 & 0 & z & 0\\0 & z & 0 & z\\z & 0 & z & 0\end{pmatrix}\]
\end{minipage}%
\par\noindent
\begin{minipage}[c]{0.495\linewidth}\centering\fontsize{9.7}{12}\selectfont
\setlength{\abovedisplayskip}{4pt}\setlength{\belowdisplayskip}{4pt}
\setlength{\abovedisplayshortskip}{4pt}\setlength{\belowdisplayshortskip}{4pt}
\setlength{\arraycolsep}{3pt}\renewcommand{\arraystretch}{1.08}
\[A_{+}(1)=\begin{pmatrix}2 & 0 & 0 & 0\\0 & 1 & 0 & 0\\0 & 0 & 1 & 0\\0 & 0 & 0 & 1\end{pmatrix}\]
\end{minipage}%
\begin{minipage}[c]{0.495\linewidth}\centering\fontsize{9.7}{12}\selectfont
\setlength{\abovedisplayskip}{4pt}\setlength{\belowdisplayskip}{4pt}
\setlength{\abovedisplayshortskip}{4pt}\setlength{\belowdisplayshortskip}{4pt}
\setlength{\arraycolsep}{3pt}\renewcommand{\arraystretch}{1.08}
\[A_{-}(1)=\begin{pmatrix}2 & 0 & 0 & -2\\0 & 2 & -1 & 0\\0 & -1 & 2 & -1\\-1 & 0 & -1 & 2\end{pmatrix}\]
\end{minipage}%
\par\smallskip\noindent\small $h_+=3,\quad h_-=9,\quad \Omega=12.$\quad Vertices per slice: 9.
\end{minipage}\par\vfill
\noindent\begin{minipage}{\linewidth}
\noindent\textbf{Class 28}\quad\small (constant ID 1493)\hfill $\symbf{r}=(2,2,2,2)$
\par\smallskip\noindent\textcolor{rulegray}{\rule{\linewidth}{0.3pt}}
\begin{minipage}[c]{0.495\linewidth}\centering\fontsize{9.7}{12}\selectfont
\setlength{\abovedisplayskip}{4pt}\setlength{\belowdisplayskip}{4pt}
\setlength{\abovedisplayshortskip}{4pt}\setlength{\belowdisplayshortskip}{4pt}
\setlength{\arraycolsep}{3pt}\renewcommand{\arraystretch}{1.08}
\[N_{+}^{(28)}(z)=\begin{pmatrix}0 & 0 & 0 & z\\0 & 0 & 0 & z\\0 & 0 & 0 & z\\z & z & z & 0\end{pmatrix}\]
\end{minipage}%
\begin{minipage}[c]{0.495\linewidth}\centering\fontsize{9.7}{12}\selectfont
\setlength{\abovedisplayskip}{4pt}\setlength{\belowdisplayskip}{4pt}
\setlength{\abovedisplayshortskip}{4pt}\setlength{\belowdisplayshortskip}{4pt}
\setlength{\arraycolsep}{3pt}\renewcommand{\arraystretch}{1.08}
\[N_{-}^{(28)}(z)=\begin{pmatrix}0 & 0 & z & 0\\0 & z & 0 & 0\\z & 0 & 0 & 0\\0 & 0 & 0 & z\end{pmatrix}\]
\end{minipage}%
\par\noindent
\begin{minipage}[c]{0.495\linewidth}\centering\fontsize{9.7}{12}\selectfont
\setlength{\abovedisplayskip}{4pt}\setlength{\belowdisplayskip}{4pt}
\setlength{\abovedisplayshortskip}{4pt}\setlength{\belowdisplayshortskip}{4pt}
\setlength{\arraycolsep}{3pt}\renewcommand{\arraystretch}{1.08}
\[A_{+}(1)=\begin{pmatrix}2 & 0 & 0 & -1\\0 & 2 & 0 & -1\\0 & 0 & 2 & -1\\-1 & -1 & -1 & 2\end{pmatrix}\]
\end{minipage}%
\begin{minipage}[c]{0.495\linewidth}\centering\fontsize{9.7}{12}\selectfont
\setlength{\abovedisplayskip}{4pt}\setlength{\belowdisplayskip}{4pt}
\setlength{\abovedisplayshortskip}{4pt}\setlength{\belowdisplayshortskip}{4pt}
\setlength{\arraycolsep}{3pt}\renewcommand{\arraystretch}{1.08}
\[A_{-}(1)=\begin{pmatrix}2 & 0 & -1 & 0\\0 & 1 & 0 & 0\\-1 & 0 & 2 & 0\\0 & 0 & 0 & 1\end{pmatrix}\]
\end{minipage}%
\par\smallskip\noindent\small $h_+=6,\quad h_-=3,\quad \Omega=18.$\quad Vertices per slice: 8.
\end{minipage}\par\vfill
\noindent\begin{minipage}{\linewidth}
\noindent\textbf{Class 29}\quad\small (constant ID 1494)\hfill $\symbf{r}=(2,2,2,2)$
\par\smallskip\noindent\textcolor{rulegray}{\rule{\linewidth}{0.3pt}}
\begin{minipage}[c]{0.495\linewidth}\centering\fontsize{9.7}{12}\selectfont
\setlength{\abovedisplayskip}{4pt}\setlength{\belowdisplayskip}{4pt}
\setlength{\abovedisplayshortskip}{4pt}\setlength{\belowdisplayshortskip}{4pt}
\setlength{\arraycolsep}{3pt}\renewcommand{\arraystretch}{1.08}
\[N_{+}^{(29)}(z)=\begin{pmatrix}0 & 0 & 0 & z\\0 & 0 & 0 & z\\0 & 0 & 0 & z\\z & z & z & 0\end{pmatrix}\]
\end{minipage}%
\begin{minipage}[c]{0.495\linewidth}\centering\fontsize{9.7}{12}\selectfont
\setlength{\abovedisplayskip}{4pt}\setlength{\belowdisplayskip}{4pt}
\setlength{\abovedisplayshortskip}{4pt}\setlength{\belowdisplayshortskip}{4pt}
\setlength{\arraycolsep}{3pt}\renewcommand{\arraystretch}{1.08}
\[N_{-}^{(29)}(z)=\begin{pmatrix}z & 0 & 0 & 0\\0 & z & 0 & 0\\0 & 0 & z & 0\\0 & 0 & 0 & z\end{pmatrix}\]
\end{minipage}%
\par\noindent
\begin{minipage}[c]{0.495\linewidth}\centering\fontsize{9.7}{12}\selectfont
\setlength{\abovedisplayskip}{4pt}\setlength{\belowdisplayskip}{4pt}
\setlength{\abovedisplayshortskip}{4pt}\setlength{\belowdisplayshortskip}{4pt}
\setlength{\arraycolsep}{3pt}\renewcommand{\arraystretch}{1.08}
\[A_{+}(1)=\begin{pmatrix}2 & 0 & 0 & -1\\0 & 2 & 0 & -1\\0 & 0 & 2 & -1\\-1 & -1 & -1 & 2\end{pmatrix}\]
\end{minipage}%
\begin{minipage}[c]{0.495\linewidth}\centering\fontsize{9.7}{12}\selectfont
\setlength{\abovedisplayskip}{4pt}\setlength{\belowdisplayskip}{4pt}
\setlength{\abovedisplayshortskip}{4pt}\setlength{\belowdisplayshortskip}{4pt}
\setlength{\arraycolsep}{3pt}\renewcommand{\arraystretch}{1.08}
\[A_{-}(1)=\begin{pmatrix}1 & 0 & 0 & 0\\0 & 1 & 0 & 0\\0 & 0 & 1 & 0\\0 & 0 & 0 & 1\end{pmatrix}\]
\end{minipage}%
\par\smallskip\noindent\small $h_+=6,\quad h_-=3,\quad \Omega=9.$\quad Vertices per slice: 8.
\end{minipage}\par\vfill
\clearpage
\subsection*{Classes 30--32}
\noindent $B_\pm^{(k)}(z)=\operatorname{diag}(1+z^{r_1},\ldots,1+z^{r_4})-N_\pm^{(k)}(z)$.\par\medskip
\noindent\begin{minipage}{\linewidth}
\noindent\textbf{Class 30}\quad\small (constant ID 1834)\hfill $\symbf{r}=(2,2,2,2)$
\par\smallskip\noindent\textcolor{rulegray}{\rule{\linewidth}{0.3pt}}
\begin{minipage}[c]{0.495\linewidth}\centering\fontsize{9.7}{12}\selectfont
\setlength{\abovedisplayskip}{4pt}\setlength{\belowdisplayskip}{4pt}
\setlength{\abovedisplayshortskip}{4pt}\setlength{\belowdisplayshortskip}{4pt}
\setlength{\arraycolsep}{3pt}\renewcommand{\arraystretch}{1.08}
\[N_{+}^{(30)}(z)=\begin{pmatrix}0 & 0 & 0 & z\\0 & 0 & z & 0\\0 & z & 0 & 0\\z & 0 & 0 & 0\end{pmatrix}\]
\end{minipage}%
\begin{minipage}[c]{0.495\linewidth}\centering\fontsize{9.7}{12}\selectfont
\setlength{\abovedisplayskip}{4pt}\setlength{\belowdisplayskip}{4pt}
\setlength{\abovedisplayshortskip}{4pt}\setlength{\belowdisplayshortskip}{4pt}
\setlength{\arraycolsep}{3pt}\renewcommand{\arraystretch}{1.08}
\[N_{-}^{(30)}(z)=\begin{pmatrix}0 & 0 & z & 0\\0 & 0 & 0 & z\\z & 0 & 0 & 0\\0 & z & 0 & 0\end{pmatrix}\]
\end{minipage}%
\par\noindent
\begin{minipage}[c]{0.495\linewidth}\centering\fontsize{9.7}{12}\selectfont
\setlength{\abovedisplayskip}{4pt}\setlength{\belowdisplayskip}{4pt}
\setlength{\abovedisplayshortskip}{4pt}\setlength{\belowdisplayshortskip}{4pt}
\setlength{\arraycolsep}{3pt}\renewcommand{\arraystretch}{1.08}
\[A_{+}(1)=\begin{pmatrix}2 & 0 & 0 & -1\\0 & 2 & -1 & 0\\0 & -1 & 2 & 0\\-1 & 0 & 0 & 2\end{pmatrix}\]
\end{minipage}%
\begin{minipage}[c]{0.495\linewidth}\centering\fontsize{9.7}{12}\selectfont
\setlength{\abovedisplayskip}{4pt}\setlength{\belowdisplayskip}{4pt}
\setlength{\abovedisplayshortskip}{4pt}\setlength{\belowdisplayshortskip}{4pt}
\setlength{\arraycolsep}{3pt}\renewcommand{\arraystretch}{1.08}
\[A_{-}(1)=\begin{pmatrix}2 & 0 & -1 & 0\\0 & 2 & 0 & -1\\-1 & 0 & 2 & 0\\0 & -1 & 0 & 2\end{pmatrix}\]
\end{minipage}%
\par\smallskip\noindent\small $h_+=3,\quad h_-=3,\quad \Omega=12.$\quad Vertices per slice: 4.
\end{minipage}\par\vfill
\noindent\begin{minipage}{\linewidth}
\noindent\textbf{Class 31}\quad\small (constant ID 1836)\hfill $\symbf{r}=(2,2,2,2)$
\par\smallskip\noindent\textcolor{rulegray}{\rule{\linewidth}{0.3pt}}
\begin{minipage}[c]{0.495\linewidth}\centering\fontsize{9.7}{12}\selectfont
\setlength{\abovedisplayskip}{4pt}\setlength{\belowdisplayskip}{4pt}
\setlength{\abovedisplayshortskip}{4pt}\setlength{\belowdisplayshortskip}{4pt}
\setlength{\arraycolsep}{3pt}\renewcommand{\arraystretch}{1.08}
\[N_{+}^{(31)}(z)=\begin{pmatrix}0 & 0 & 0 & z\\0 & 0 & z & 0\\0 & z & 0 & 0\\z & 0 & 0 & 0\end{pmatrix}\]
\end{minipage}%
\begin{minipage}[c]{0.495\linewidth}\centering\fontsize{9.7}{12}\selectfont
\setlength{\abovedisplayskip}{4pt}\setlength{\belowdisplayskip}{4pt}
\setlength{\abovedisplayshortskip}{4pt}\setlength{\belowdisplayshortskip}{4pt}
\setlength{\arraycolsep}{3pt}\renewcommand{\arraystretch}{1.08}
\[N_{-}^{(31)}(z)=\begin{pmatrix}0 & 0 & z & 0\\0 & z & 0 & z\\z & 0 & z & 0\\0 & z & 0 & 0\end{pmatrix}\]
\end{minipage}%
\par\noindent
\begin{minipage}[c]{0.495\linewidth}\centering\fontsize{9.7}{12}\selectfont
\setlength{\abovedisplayskip}{4pt}\setlength{\belowdisplayskip}{4pt}
\setlength{\abovedisplayshortskip}{4pt}\setlength{\belowdisplayshortskip}{4pt}
\setlength{\arraycolsep}{3pt}\renewcommand{\arraystretch}{1.08}
\[A_{+}(1)=\begin{pmatrix}2 & 0 & 0 & -1\\0 & 2 & -1 & 0\\0 & -1 & 2 & 0\\-1 & 0 & 0 & 2\end{pmatrix}\]
\end{minipage}%
\begin{minipage}[c]{0.495\linewidth}\centering\fontsize{9.7}{12}\selectfont
\setlength{\abovedisplayskip}{4pt}\setlength{\belowdisplayskip}{4pt}
\setlength{\abovedisplayshortskip}{4pt}\setlength{\belowdisplayshortskip}{4pt}
\setlength{\arraycolsep}{3pt}\renewcommand{\arraystretch}{1.08}
\[A_{-}(1)=\begin{pmatrix}2 & 0 & -1 & 0\\0 & 1 & 0 & -1\\-1 & 0 & 1 & 0\\0 & -1 & 0 & 2\end{pmatrix}\]
\end{minipage}%
\par\smallskip\noindent\small $h_+=3,\quad h_-=5,\quad \Omega=16.$\quad Vertices per slice: 8.
\end{minipage}\par\vfill
\noindent\begin{minipage}{\linewidth}
\noindent\textbf{Class 32}\quad\small (constant ID 1844)\hfill $\symbf{r}=(2,2,2,6)$
\par\smallskip\noindent\textcolor{rulegray}{\rule{\linewidth}{0.3pt}}
\begin{minipage}[c]{0.495\linewidth}\centering\fontsize{9.7}{12}\selectfont
\setlength{\abovedisplayskip}{4pt}\setlength{\belowdisplayskip}{4pt}
\setlength{\abovedisplayshortskip}{4pt}\setlength{\belowdisplayshortskip}{4pt}
\setlength{\arraycolsep}{3pt}\renewcommand{\arraystretch}{1.08}
\[N_{+}^{(32)}(z)=\begin{pmatrix}0 & 0 & 0 & z\\0 & 0 & z & 0\\0 & z & 0 & 0\\z^{5} + z & 0 & z^{3} & 0\end{pmatrix}\]
\end{minipage}%
\begin{minipage}[c]{0.495\linewidth}\centering\fontsize{9.7}{12}\selectfont
\setlength{\abovedisplayskip}{4pt}\setlength{\belowdisplayskip}{4pt}
\setlength{\abovedisplayshortskip}{4pt}\setlength{\belowdisplayshortskip}{4pt}
\setlength{\arraycolsep}{3pt}\renewcommand{\arraystretch}{1.08}
\[N_{-}^{(32)}(z)=\begin{pmatrix}0 & z & 0 & 0\\z & 0 & 0 & 0\\0 & 0 & 0 & z\\z^{3} & 0 & z^{5} + z & 0\end{pmatrix}\]
\end{minipage}%
\par\noindent
\begin{minipage}[c]{0.495\linewidth}\centering\fontsize{9.7}{12}\selectfont
\setlength{\abovedisplayskip}{4pt}\setlength{\belowdisplayskip}{4pt}
\setlength{\abovedisplayshortskip}{4pt}\setlength{\belowdisplayshortskip}{4pt}
\setlength{\arraycolsep}{3pt}\renewcommand{\arraystretch}{1.08}
\[A_{+}(1)=\begin{pmatrix}2 & 0 & 0 & -1\\0 & 2 & -1 & 0\\0 & -1 & 2 & 0\\-2 & 0 & -1 & 2\end{pmatrix}\]
\end{minipage}%
\begin{minipage}[c]{0.495\linewidth}\centering\fontsize{9.7}{12}\selectfont
\setlength{\abovedisplayskip}{4pt}\setlength{\belowdisplayskip}{4pt}
\setlength{\abovedisplayshortskip}{4pt}\setlength{\belowdisplayshortskip}{4pt}
\setlength{\arraycolsep}{3pt}\renewcommand{\arraystretch}{1.08}
\[A_{-}(1)=\begin{pmatrix}2 & -1 & 0 & 0\\-1 & 2 & 0 & 0\\0 & 0 & 2 & -1\\-1 & 0 & -2 & 2\end{pmatrix}\]
\end{minipage}%
\par\smallskip\noindent\small $h_+=8,\quad h_-=8,\quad \Omega=16.$\quad Vertices per slice: 6.
\end{minipage}\par\vfill
\clearpage
\subsection*{Classes 33--35}
\noindent $B_\pm^{(k)}(z)=\operatorname{diag}(1+z^{r_1},\ldots,1+z^{r_4})-N_\pm^{(k)}(z)$.\par\medskip
\noindent\begin{minipage}{\linewidth}
\noindent\textbf{Class 33}\quad\small (constant ID 1848)\hfill $\symbf{r}=(2,2,2,10)$
\par\smallskip\noindent\textcolor{rulegray}{\rule{\linewidth}{0.3pt}}
\begin{minipage}[c]{0.495\linewidth}\centering\fontsize{9.7}{12}\selectfont
\setlength{\abovedisplayskip}{4pt}\setlength{\belowdisplayskip}{4pt}
\setlength{\abovedisplayshortskip}{4pt}\setlength{\belowdisplayshortskip}{4pt}
\setlength{\arraycolsep}{3pt}\renewcommand{\arraystretch}{1.08}
\[N_{+}^{(33)}(z)=\begin{pmatrix}0 & 0 & 0 & z\\0 & 0 & z & 0\\0 & z & 0 & 0\\z^{9} + z^{5} + z & 0 & z^{7} + z^{3} & 0\end{pmatrix}\]
\end{minipage}%
\begin{minipage}[c]{0.495\linewidth}\centering\fontsize{9.7}{12}\selectfont
\setlength{\abovedisplayskip}{4pt}\setlength{\belowdisplayskip}{4pt}
\setlength{\abovedisplayshortskip}{4pt}\setlength{\belowdisplayshortskip}{4pt}
\setlength{\arraycolsep}{3pt}\renewcommand{\arraystretch}{1.08}
\[N_{-}^{(33)}(z)=\begin{pmatrix}0 & z & 0 & 0\\z & 0 & 0 & 0\\0 & 0 & 0 & z\\z^{7} + z^{3} & 0 & z^{9} + z^{5} + z & 0\end{pmatrix}\]
\end{minipage}%
\par\noindent
\begin{minipage}[c]{0.495\linewidth}\centering\fontsize{9.7}{12}\selectfont
\setlength{\abovedisplayskip}{4pt}\setlength{\belowdisplayskip}{4pt}
\setlength{\abovedisplayshortskip}{4pt}\setlength{\belowdisplayshortskip}{4pt}
\setlength{\arraycolsep}{3pt}\renewcommand{\arraystretch}{1.08}
\[A_{+}(1)=\begin{pmatrix}2 & 0 & 0 & -1\\0 & 2 & -1 & 0\\0 & -1 & 2 & 0\\-3 & 0 & -2 & 2\end{pmatrix}\]
\end{minipage}%
\begin{minipage}[c]{0.495\linewidth}\centering\fontsize{9.7}{12}\selectfont
\setlength{\abovedisplayskip}{4pt}\setlength{\belowdisplayskip}{4pt}
\setlength{\abovedisplayshortskip}{4pt}\setlength{\belowdisplayshortskip}{4pt}
\setlength{\arraycolsep}{3pt}\renewcommand{\arraystretch}{1.08}
\[A_{-}(1)=\begin{pmatrix}2 & -1 & 0 & 0\\-1 & 2 & 0 & 0\\0 & 0 & 2 & -1\\-2 & 0 & -3 & 2\end{pmatrix}\]
\end{minipage}%
\par\smallskip\noindent\small $h_+=18,\quad h_-=18,\quad \Omega=36.$\quad Vertices per slice: 8.
\end{minipage}\par\vfill
\noindent\begin{minipage}{\linewidth}
\noindent\textbf{Class 34}\quad\small (constant ID 1852)\hfill $\symbf{r}=(2,2,2,2)$
\par\smallskip\noindent\textcolor{rulegray}{\rule{\linewidth}{0.3pt}}
\begin{minipage}[c]{0.495\linewidth}\centering\fontsize{9.7}{12}\selectfont
\setlength{\abovedisplayskip}{4pt}\setlength{\belowdisplayskip}{4pt}
\setlength{\abovedisplayshortskip}{4pt}\setlength{\belowdisplayshortskip}{4pt}
\setlength{\arraycolsep}{3pt}\renewcommand{\arraystretch}{1.08}
\[N_{+}^{(34)}(z)=\begin{pmatrix}0 & 0 & 0 & z\\0 & 0 & z & 0\\0 & z & 0 & z\\z & 0 & z & 0\end{pmatrix}\]
\end{minipage}%
\begin{minipage}[c]{0.495\linewidth}\centering\fontsize{9.7}{12}\selectfont
\setlength{\abovedisplayskip}{4pt}\setlength{\belowdisplayskip}{4pt}
\setlength{\abovedisplayshortskip}{4pt}\setlength{\belowdisplayshortskip}{4pt}
\setlength{\arraycolsep}{3pt}\renewcommand{\arraystretch}{1.08}
\[N_{-}^{(34)}(z)=\begin{pmatrix}0 & 0 & z & 0\\0 & 0 & 0 & z\\z & 0 & z & 0\\0 & z & 0 & z\end{pmatrix}\]
\end{minipage}%
\par\noindent
\begin{minipage}[c]{0.495\linewidth}\centering\fontsize{9.7}{12}\selectfont
\setlength{\abovedisplayskip}{4pt}\setlength{\belowdisplayskip}{4pt}
\setlength{\abovedisplayshortskip}{4pt}\setlength{\belowdisplayshortskip}{4pt}
\setlength{\arraycolsep}{3pt}\renewcommand{\arraystretch}{1.08}
\[A_{+}(1)=\begin{pmatrix}2 & 0 & 0 & -1\\0 & 2 & -1 & 0\\0 & -1 & 2 & -1\\-1 & 0 & -1 & 2\end{pmatrix}\]
\end{minipage}%
\begin{minipage}[c]{0.495\linewidth}\centering\fontsize{9.7}{12}\selectfont
\setlength{\abovedisplayskip}{4pt}\setlength{\belowdisplayskip}{4pt}
\setlength{\abovedisplayshortskip}{4pt}\setlength{\belowdisplayshortskip}{4pt}
\setlength{\arraycolsep}{3pt}\renewcommand{\arraystretch}{1.08}
\[A_{-}(1)=\begin{pmatrix}2 & 0 & -1 & 0\\0 & 2 & 0 & -1\\-1 & 0 & 1 & 0\\0 & -1 & 0 & 1\end{pmatrix}\]
\end{minipage}%
\par\smallskip\noindent\small $h_+=5,\quad h_-=5,\quad \Omega=20.$\quad Vertices per slice: 8.
\end{minipage}\par\vfill
\noindent\begin{minipage}{\linewidth}
\noindent\textbf{Class 35}\quad\small (constant ID 1853)\hfill $\symbf{r}=(2,2,2,2)$
\par\smallskip\noindent\textcolor{rulegray}{\rule{\linewidth}{0.3pt}}
\begin{minipage}[c]{0.495\linewidth}\centering\fontsize{9.7}{12}\selectfont
\setlength{\abovedisplayskip}{4pt}\setlength{\belowdisplayskip}{4pt}
\setlength{\abovedisplayshortskip}{4pt}\setlength{\belowdisplayshortskip}{4pt}
\setlength{\arraycolsep}{3pt}\renewcommand{\arraystretch}{1.08}
\[N_{+}^{(35)}(z)=\begin{pmatrix}0 & 0 & 0 & z\\0 & 0 & z & 0\\0 & z & 0 & z\\z & 0 & z & 0\end{pmatrix}\]
\end{minipage}%
\begin{minipage}[c]{0.495\linewidth}\centering\fontsize{9.7}{12}\selectfont
\setlength{\abovedisplayskip}{4pt}\setlength{\belowdisplayskip}{4pt}
\setlength{\abovedisplayshortskip}{4pt}\setlength{\belowdisplayshortskip}{4pt}
\setlength{\arraycolsep}{3pt}\renewcommand{\arraystretch}{1.08}
\[N_{-}^{(35)}(z)=\begin{pmatrix}z & 0 & 0 & 0\\0 & z & 0 & 0\\0 & 0 & z & 0\\0 & 0 & 0 & z\end{pmatrix}\]
\end{minipage}%
\par\noindent
\begin{minipage}[c]{0.495\linewidth}\centering\fontsize{9.7}{12}\selectfont
\setlength{\abovedisplayskip}{4pt}\setlength{\belowdisplayskip}{4pt}
\setlength{\abovedisplayshortskip}{4pt}\setlength{\belowdisplayshortskip}{4pt}
\setlength{\arraycolsep}{3pt}\renewcommand{\arraystretch}{1.08}
\[A_{+}(1)=\begin{pmatrix}2 & 0 & 0 & -1\\0 & 2 & -1 & 0\\0 & -1 & 2 & -1\\-1 & 0 & -1 & 2\end{pmatrix}\]
\end{minipage}%
\begin{minipage}[c]{0.495\linewidth}\centering\fontsize{9.7}{12}\selectfont
\setlength{\abovedisplayskip}{4pt}\setlength{\belowdisplayskip}{4pt}
\setlength{\abovedisplayshortskip}{4pt}\setlength{\belowdisplayshortskip}{4pt}
\setlength{\arraycolsep}{3pt}\renewcommand{\arraystretch}{1.08}
\[A_{-}(1)=\begin{pmatrix}1 & 0 & 0 & 0\\0 & 1 & 0 & 0\\0 & 0 & 1 & 0\\0 & 0 & 0 & 1\end{pmatrix}\]
\end{minipage}%
\par\smallskip\noindent\small $h_+=5,\quad h_-=3,\quad \Omega=16.$\quad Vertices per slice: 8.
\end{minipage}\par\vfill
\clearpage
\subsection*{Classes 36--37}
\noindent $B_\pm^{(k)}(z)=\operatorname{diag}(1+z^{r_1},\ldots,1+z^{r_4})-N_\pm^{(k)}(z)$.\par\medskip
\noindent\begin{minipage}{\linewidth}
\noindent\textbf{Class 36}\quad\small (constant ID 1902)\hfill $\symbf{r}=(2,2,2,3)$
\par\smallskip\noindent\textcolor{rulegray}{\rule{\linewidth}{0.3pt}}
\begin{minipage}[c]{0.495\linewidth}\centering\fontsize{9.7}{12}\selectfont
\setlength{\abovedisplayskip}{4pt}\setlength{\belowdisplayskip}{4pt}
\setlength{\abovedisplayshortskip}{4pt}\setlength{\belowdisplayshortskip}{4pt}
\setlength{\arraycolsep}{3pt}\renewcommand{\arraystretch}{1.08}
\[N_{+}^{(36)}(z)=\begin{pmatrix}0 & 0 & 0 & z\\0 & 0 & z & 0\\0 & z & z & 0\\z^{2} + z & 0 & 0 & 0\end{pmatrix}\]
\end{minipage}%
\begin{minipage}[c]{0.495\linewidth}\centering\fontsize{9.7}{12}\selectfont
\setlength{\abovedisplayskip}{4pt}\setlength{\belowdisplayskip}{4pt}
\setlength{\abovedisplayshortskip}{4pt}\setlength{\belowdisplayshortskip}{4pt}
\setlength{\arraycolsep}{3pt}\renewcommand{\arraystretch}{1.08}
\[N_{-}^{(36)}(z)=\begin{pmatrix}z & z & 0 & 0\\z & 0 & 0 & 0\\0 & 0 & 0 & z\\0 & 0 & z^{2} + z & 0\end{pmatrix}\]
\end{minipage}%
\par\noindent
\begin{minipage}[c]{0.495\linewidth}\centering\fontsize{9.7}{12}\selectfont
\setlength{\abovedisplayskip}{4pt}\setlength{\belowdisplayskip}{4pt}
\setlength{\abovedisplayshortskip}{4pt}\setlength{\belowdisplayshortskip}{4pt}
\setlength{\arraycolsep}{3pt}\renewcommand{\arraystretch}{1.08}
\[A_{+}(1)=\begin{pmatrix}2 & 0 & 0 & -1\\0 & 2 & -1 & 0\\0 & -1 & 1 & 0\\-2 & 0 & 0 & 2\end{pmatrix}\]
\end{minipage}%
\begin{minipage}[c]{0.495\linewidth}\centering\fontsize{9.7}{12}\selectfont
\setlength{\abovedisplayskip}{4pt}\setlength{\belowdisplayskip}{4pt}
\setlength{\abovedisplayshortskip}{4pt}\setlength{\belowdisplayshortskip}{4pt}
\setlength{\arraycolsep}{3pt}\renewcommand{\arraystretch}{1.08}
\[A_{-}(1)=\begin{pmatrix}1 & -1 & 0 & 0\\-1 & 2 & 0 & 0\\0 & 0 & 2 & -1\\0 & 0 & -2 & 2\end{pmatrix}\]
\end{minipage}%
\par\smallskip\noindent\small $h_+=5,\quad h_-=5,\quad \Omega=10.$\quad Vertices per slice: 9.
\end{minipage}\par\vfill
\noindent\begin{minipage}{\linewidth}
\noindent\textbf{Class 37}\quad\small (constant ID 2614)\hfill $\symbf{r}=(2,2,6,6)$
\par\smallskip\noindent\textcolor{rulegray}{\rule{\linewidth}{0.3pt}}
\begin{minipage}[c]{0.495\linewidth}\centering\fontsize{9.7}{12}\selectfont
\setlength{\abovedisplayskip}{4pt}\setlength{\belowdisplayskip}{4pt}
\setlength{\abovedisplayshortskip}{4pt}\setlength{\belowdisplayshortskip}{4pt}
\setlength{\arraycolsep}{3pt}\renewcommand{\arraystretch}{1.08}
\[N_{+}^{(37)}(z)=\begin{pmatrix}0 & 0 & 0 & z\\0 & 0 & z & 0\\z^{3} & z^{5} + z & 0 & 0\\z^{5} + z & z^{3} & 0 & 0\end{pmatrix}\]
\end{minipage}%
\begin{minipage}[c]{0.495\linewidth}\centering\fontsize{9.7}{12}\selectfont
\setlength{\abovedisplayskip}{4pt}\setlength{\belowdisplayskip}{4pt}
\setlength{\abovedisplayshortskip}{4pt}\setlength{\belowdisplayshortskip}{4pt}
\setlength{\arraycolsep}{3pt}\renewcommand{\arraystretch}{1.08}
\[N_{-}^{(37)}(z)=\begin{pmatrix}0 & 0 & z & 0\\0 & 0 & 0 & z\\z^{5} + z & z^{3} & 0 & 0\\z^{3} & z^{5} + z & 0 & 0\end{pmatrix}\]
\end{minipage}%
\par\noindent
\begin{minipage}[c]{0.495\linewidth}\centering\fontsize{9.7}{12}\selectfont
\setlength{\abovedisplayskip}{4pt}\setlength{\belowdisplayskip}{4pt}
\setlength{\abovedisplayshortskip}{4pt}\setlength{\belowdisplayshortskip}{4pt}
\setlength{\arraycolsep}{3pt}\renewcommand{\arraystretch}{1.08}
\[A_{+}(1)=\begin{pmatrix}2 & 0 & 0 & -1\\0 & 2 & -1 & 0\\-1 & -2 & 2 & 0\\-2 & -1 & 0 & 2\end{pmatrix}\]
\end{minipage}%
\begin{minipage}[c]{0.495\linewidth}\centering\fontsize{9.7}{12}\selectfont
\setlength{\abovedisplayskip}{4pt}\setlength{\belowdisplayskip}{4pt}
\setlength{\abovedisplayshortskip}{4pt}\setlength{\belowdisplayshortskip}{4pt}
\setlength{\arraycolsep}{3pt}\renewcommand{\arraystretch}{1.08}
\[A_{-}(1)=\begin{pmatrix}2 & 0 & -1 & 0\\0 & 2 & 0 & -1\\-2 & -1 & 2 & 0\\-1 & -2 & 0 & 2\end{pmatrix}\]
\end{minipage}%
\par\smallskip\noindent\small $h_+=12,\quad h_-=12,\quad \Omega=24.$\quad Vertices per slice: 8.
\end{minipage}\par\vfill


\clearpage
\appendix
\section{Reproducibility and proof dependencies}
This classification uses exact arithmetic and symbolic feasibility checks.
The finite-type implication from simultaneous positivity, the coefficient
bound, the enumeration scheme and the chain obstruction are proved above.
The finite enumeration and the remaining feasibility and coverage statements
are supported by executable exact-arithmetic checks. The computation trusts
the C++ and Python implementations, SymPy rational arithmetic and Z3's UNSAT
answers; it is not a proof-assistant formalization.

The checked runtime is Python 3.10.6, SymPy 1.14.0 and Z3 5.1.0. The
source files are under \code{research/rank4/}. The mutation routines and
installed Z3 package are reused from \code{research/rank3/}.

\begin{center}\small\renewcommand{\arraystretch}{1.15}
\begin{tabular}{@{}p{0.44\linewidth}p{0.51\linewidth}@{}}
\toprule
File & Content\\\midrule
\code{enumerate_constants.cpp} & Exhaustive ordered-vector enumeration.\\
\code{constant_candidates.json} & The $4,865$ canonical constant pairs.\\
\code{constant_verification.json} & Integer positive vectors and parity exclusions.\\
\code{lift_feasibility.jsonl} & All lift decisions and the $37$ witnesses.\\
\code{families.jsonl} & Linear equations and exact periodicity certificates.\\
\code{smt_queries.zip} & Complete exclusion and coverage formulas.\\
\code{verification.json} & Replayed answers, encodings, versions and hashes.\\
\code{slice_signatures.json} & Full canonical slice encodings.\\
\code{classification.json} & The catalogue, including both matrices at $1$.\\
\bottomrule
\end{tabular}
\end{center}

Run the commands in \code{REPRODUCE.md} from the repository root. The
enumerator has no delay parameter. The lift solver preserves UNKNOWN
answers and the verifier refuses a complete classification while any remain.
The final replay contains $4,865$ UNSAT answers: $4,407$ have elementary
parity or chain exclusions, $421$ exclude the remaining impossible lifts,
and $37$ exclude lifts outside the displayed families.

\paragraph{Building this document.}
Run \code{python research/rank4/make_pdf_tables.py}, then build
\code{rank4-classification.tex} with LuaLaTeX and \code{latexmk}, writing
to \code{output/pdf/}. The generator evaluates every polynomial matrix at
$1$ and checks it against the separately enumerated constants.

\begin{thebibliography}{9}
\bibitem{Mizuno}
Y.~Mizuno, \emph{Periodic Y-systems and Nahm sums: the rank 2 case},
SIGMA \textbf{21} (2025), 094, 17 pp.
\href{https://doi.org/10.3842/SIGMA.2025.094}{doi:10.3842/SIGMA.2025.094}.
\href{https://arxiv.org/abs/2301.13239v2}{arXiv:2301.13239v2}.
\bibitem{TData}
Y.~Mizuno, \emph{Difference equations arising from cluster algebras},
\href{https://arxiv.org/abs/1912.05710v2}{arXiv:1912.05710v2} (2020).
\end{thebibliography}
\end{document}
